% Transcription of Cayley's Collected Papers, Volume 11, pages 501--550
% Source: collectedmathema11cayluoft.pdf
% Articles: 786 (Equation, cont.), 787 (Function), 788 (Galois), 789 (Gauss), 790 (Geometry)
\documentclass[11pt,a4paper]{article}
\usepackage[utf8]{inputenc}
\usepackage[T1]{fontenc}
\usepackage[english]{babel}
\usepackage{amsmath,amssymb,amsthm}
\usepackage{geometry}
\usepackage{titlesec}
\usepackage{fancyhdr}
\usepackage{microtype}

\geometry{margin=2.5cm}

\pagestyle{fancy}
\fancyhf{}
\renewcommand{\headrulewidth}{0pt}
\fancyhead[C]{\thepage}

\newcommand{\E}{\mathrm{e}}
\DeclareMathOperator{\am}{am}
\DeclareMathOperator{\sn}{sn}
\DeclareMathOperator{\cn}{cn}
\DeclareMathOperator{\dn}{dn}

\titleformat{\section}{\large\bfseries\centering}{}{0em}{}
\titleformat{\subsection}{\normalsize\bfseries}{}{0em}{}
\titlespacing*{\section}{0pt}{\baselineskip}{\baselineskip}
\titlespacing*{\subsection}{0pt}{\baselineskip}{0.5\baselineskip}

\begin{document}
\thispagestyle{empty}
\begin{center}
\Large\textbf{The Collected Mathematical Papers of Arthur Cayley}\\[6pt]
\large Volume XI\\[12pt]
\normalsize Pages 501--550\\[6pt]
\small Transcribed from \textit{Collected Mathematical Papers, Vol.~11} (Cambridge, 1896)
\end{center}
\vspace{2\baselineskip}

% ============================================================
\section*[Article 786]{\S\,786. EQUATION (continued).}
\addcontentsline{toc}{section}{786. Equation (continued)}
% ============================================================

\noindent[From the \textit{Encyclop{\ae}dia Britannica}, Ninth Edition, vol.~viii (1878), pp.~610--627.]

\subsection*{\S\,15.}

It will be recollected that the expression \textit{number} and the correlative epithet \textit{numerical} were at the outset used in a wide sense, as extending to imaginaries. This extension arises out of the theory of equations by a process analogous to that by which number, in its original most restricted sense of positive integer number, was extended to have the meaning of a real positive or negative magnitude susceptible of continuous variation.

If for a moment number is understood in its most restricted sense as meaning positive integer number, the solution of a simple equation leads to an extension; $ax - b = 0$, gives $x = \frac{b}{a}$, a positive fraction, and we can in this manner represent, not accurately, but as nearly as we please, any positive magnitude whatever; so an equation $ax + b = 0$ gives $x = -\frac{b}{a}$, which (approximately as before) represents any negative magnitude. We thus arrive at the extended signification of number as a continuously varying positive or negative magnitude. Such numbers may be added or subtracted, multiplied or divided one by another, and the result is always a number. Now from a quadric equation we derive, in like manner, the notion of a complex or imaginary number such as is spoken of above. The equation $x^{2} + 1 = 0$ is not (in the foregoing sense, number = real number) satisfied by any numerical value whatever of $x$; but we assume that there is a number which we call $i$, satisfying the equation $i^{2} + 1 = 0$; and then taking $a$ and $b$ any real numbers, we form an expression such as $a + bi$, and use the expression number in this extended sense: any two such numbers may be added or subtracted, multiplied or divided one by the other, and the result is always a number. And if we consider first a quadric equation $x^{2} + px + q = 0$ where $p$ and $q$ are real numbers, and next the like equation, where $p$ and $q$ are any numbers whatever, it can be shown that there exists for $x$ a numerical value which satisfies the equation; or, in other words, it can be shown that the equation has a numerical root. The like theorem, in fact, holds good for an equation of any order whatever. But suppose for a moment that this was not the case: say that there was a cubic equation $x^{3} + px^{2} + qx + r = 0$, with numerical coefficients, not satisfied by any numerical value of $x$, we should have to establish a new imaginary $j$ satisfying some such equation, and should then have to consider numbers of the form $a + bj$, or perhaps $a + bj + cj^{2}$ ($a$, $b$, $c$ numbers $a + \beta i$ of the kind heretofore considered); first we should be thrown back on the quadric equation $x^{2} + px + q = 0$, $p$ and $q$ being now numbers of the last-mentioned extended form---non constat that every such equation has a numerical root---and if not, we might be led to other imaginaries $k$, $l$, \&c., and so on \textit{ad infinitum} in inextricable confusion.

But in fact a numerical equation of any order whatever has always a numerical root, and thus numbers (in the foregoing sense, number = quantity of the form $a + \beta i$) form (what real numbers do not) a universe complete in itself, such that starting in it we are never led out of it. There may very well be, and perhaps are, numbers in a more general sense of the term (quaternions are not a case in point, as the ordinary laws of combination are not adhered to): but in order to have to do with such numbers (if any), we must start with them.

\subsection*{\S\,16.}

The theorem as regards numerical equations thus is, every numerical equation has a numerical root; or for shortness (the meaning being as before), every equation has a root. Of course the theorem is the reverse of self-evident, and it requires proof; but provisionally assuming it as true, we derive from it the general theory of numerical equations. As the term root was introduced in the course of an explanation, it will be convenient to give here the formal definition.

A number $a$ such that substituted for $x$ it makes the function $x^{n} - p_{1}x^{n-1} + \ldots \pm p_{n}$ to be $=0$, or say such that it satisfies the equation $f(x) = 0$, is said to be a \textit{root} of the equation; that is, $a$ being a root, we have
\[
a^{n} - p_{1}a^{n-1} + \ldots \pm p_{n} = 0, \quad \text{or say} \quad f(a) = 0;
\]
and it is then easily shown that $x - a$ is a factor of the function $f(x)$, viz.\ that we have $f(x) = (x - a)f_{1}(x)$, where $f_{1}(x)$ is a function $x^{n-1} - q_{1}x^{n-2} + \ldots \pm q_{n-1}$ of the order $n - 1$, with numerical coefficients $q_{1}, q_{2}, \ldots, q_{n-1}$.

In general, $a$ is not a root of the equation $f_{1}(x) = 0$, but it may be so; i.e.\ $f_{1}(x)$ may contain the factor $x - a$; when this is so, $f(x)$ will contain the factor $(x - a)^{2}$; writing then $f(x) = (x - a)^{2}f_{2}(x)$, and assuming that $a$ is not a root of the equation $f_{2}(x) = 0$, $x = a$ is then said to be a \textit{double root} of the equation $f(x) = 0$; and similarly $f(x)$ may contain the factor $(x - a)^{3}$ and no higher power, and $x = a$ is then a \textit{triple root}; and so on.

Supposing, in general, that $f(x) = (x - a)^{\alpha}F(x)$, $\alpha$ being a positive integer which may be $=1$, $(x - a)^{\alpha}$ the highest power of $x - a$ which divides $f(x)$, and $F(x)$ being of course of the order $n - \alpha$, then the equation $F(x) = 0$ will have a root $b$ which will be different from $a$; $x - b$ will be a factor, in general a simple one, but it may be a multiple one, of $F(x)$, and $f(x)$ will in this case be $= (x - a)^{\alpha}(x - b)^{\beta}\Phi(x)$, $\beta$ a positive integer which may be $=1$, $(x - b)^{\beta}$ the highest power of $x - b$ in $F(x)$ or $f(x)$, and $\Phi(x)$ being of course of the order $n - \alpha - \beta$.

The original equation $f(x) = 0$ is in this case said to have $\alpha$ roots each $=a$, $\beta$ roots each $=b$; and so on for any other factors $(x - c)^{\gamma}$, \&c.

We have thus the theorem---A numerical equation of the order $n$ has in every case $n$ roots, viz.\ there exist $n$ numbers $a, b, \ldots$, in general all distinct, but which may arrange themselves in any sets of equal values, such that $f(x) = (x - a)(x - b)(x - c)\ldots$ identically.

If the equation has equal roots, these can in general be determined: and the case is at any rate a special one which may be in the first instance excluded from consideration. It is therefore, in general, assumed that the equation $f(x) = 0$ has all its roots unequal.

If the coefficients $p_{1}, p_{2}, \ldots$ are all or any one or more of them imaginary, then the equation $f(x) = 0$, separating the real and imaginary parts thereof, may be written $F(x) + i\Phi(x) = 0$, where $F(x)$, $\Phi(x)$ are each of them a function with real coefficients; and thus it appears that the equation $f(x) = 0$, with imaginary coefficients, has not in general any real root; supposing it to have a real root $a$, this must be at once a root of each of the equations $F(x) = 0$ and $\Phi(x) = 0$.

But an equation with real coefficients may have as well imaginary as real roots, and we have further the theorem that for any such equation the imaginary roots enter in pairs, viz.\ $\alpha + \beta i$ being a root, then $\alpha - \beta i$ will be also a root. It follows that, if the order be odd, there is always an odd number of real roots, and therefore at least one real root.

\subsection*{\S\,17.}

In the case of an equation with real coefficients, the question of the existence of real roots, and of their separation, has been already considered. In the general case of an equation with imaginary (it may be real) coefficients, the like question arises as to the situation of the (real or imaginary) roots; thus if, for facility of conception, we regard the constituents $\alpha$, $\beta$ of a root $\alpha + \beta i$ as the coordinates of a point in piano, and accordingly represent the root by such point, then drawing in the plane any closed curve or ``contour,'' the question is how many roots lie within such contour.

This is solved theoretically by means of a theorem of Cauchy's (1837), viz.\ writing in the original equation $x + iy$ in place of $x$, the function $f(x + iy)$ becomes $= P + iQ$, where $P$ and $Q$ are each of them a rational and integral function (with real coefficients) of $(x, y)$. Imagining the point $(x, y)$ to travel along the contour, and considering the number of changes of sign from $-$ to $+$ and from $+$ to $-$ of the fraction $\frac{P}{Q}$ corresponding to passages of the fraction through zero, that is, to values for which $P$ becomes $=0$, disregarding those for which $Q$ becomes $=0$, the difference of these numbers gives the number of roots within the contour.

It is important to remark that the demonstration does not presuppose the existence of any root; the contour may be the infinity of the plane (such infinity regarded as a contour, or closed curve), and in this case it can be shown (and that very easily) that the difference of the numbers of changes of sign is $= n$; that is, there are within the infinite contour, or (what is the same thing) there are in all, $n$ roots: thus Cauchy's theorem contains really the proof of the fundamental theorem that a numerical equation of the $n$th order (not only has a numerical root, but) has precisely $n$ roots. It would appear that this proof of the fundamental theorem in its most complete form is in principle identical with Gauss's last proof (1849) of the theorem, in the form---A numerical equation of the $n$th order has always a root\textsuperscript{*}.

\smallskip
\noindent\rule{0.4\textwidth}{0.4pt}\\
\small\textsuperscript{*} The earlier demonstrations by Euler, Lagrange, \&c., relate to the case of a numerical equation with real coefficients; and they consist in showing that such equation has always a real quadratic divisor, furnishing two roots, which are either real or else conjugate imaginaries $\alpha + \beta i$: see Lagrange's \textit{Equations Num\'{e}riques}.
\normalsize

But in the case of a finite contour, the actual determination of the difference which gives the number of real roots can be effected only in the case of a rectangular contour, by applying to each of its sides separately a method such as that of Sturm's theorem; and thus the actual determination ultimately depends on a method such as that of Sturm's theorem.

Very little has been done in regard to the calculation of the imaginary roots of an equation by approximation; and the question is not here considered.

\subsection*{\S\,18.}

A class of numerical equations which needs to be considered is that of the binomial equations $x^{n} - a = 0$ ($a = \alpha + \beta i$, a complex number). The foregoing conclusions apply, viz.\ there are always $n$ roots, which, it may be shown, are all unequal. And these can be found numerically by the extraction of the square root, and of an $n$th root, of real numbers, and by the aid of a table of natural sines and cosines\textsuperscript{*}. For writing
\[
\alpha + \beta i = \rho(\cos\lambda + i\sin\lambda),
\]
there is always a real angle $\lambda$ (positive and less than $2\pi$), such that its cosine and sine are $= \frac{\alpha}{\rho}$ and $\frac{\beta}{\rho}$ respectively; that is, writing for shortness $\sqrt{\alpha^{2} + \beta^{2}} = \rho$, we have $\alpha + \beta i = \rho(\cos\lambda + i\sin\lambda)$, or the equation is $x^{n} = \rho(\cos\lambda + i\sin\lambda)$; hence
\[
\cos\frac{n\theta}{n} + i\sin\frac{n\theta}{n} = \cos\lambda + i\sin\lambda,
\]
a value of $x$ is $= \sqrt[n]{\rho}\left(\cos\frac{\lambda}{n} + i\sin\frac{\lambda}{n}\right)$.

The formula really gives all the roots, for instead of $\lambda$ we may write $\lambda + 2s\pi$, $s$ a positive or negative integer, and then we have
\[
x = \sqrt[n]{\rho}\left(\cos\frac{\lambda + 2s\pi}{n} + i\sin\frac{\lambda + 2s\pi}{n}\right),
\]
which has the $n$ values obtained by giving to $s$ the values $0, 1, 2, \ldots, n-1$ in succession; the roots are, it is clear, represented by points lying at equal intervals on a circle. But it is more convenient to proceed somewhat differently; taking one of the roots to be $\theta$, so that $\theta^{n} = a$, then assuming $x = \theta y$, the equation becomes $y^{n} - 1 = 0$, which equation, like the original equation, has precisely $n$ roots (one of them being of course $= 1$). And the original equation $x^{n} - a = 0$ is thus reduced to the more simple equation $x^{n} - 1 = 0$; and although the theory of this equation is included in the preceding one, yet it is proper to state it separately.

The equation $x^{n} - 1 = 0$ has its several roots expressed in the form $1, \omega, \omega^{2}, \ldots, \omega^{n-1}$, where $\omega$ may be taken $= \cos\frac{2\pi}{n} + i\sin\frac{2\pi}{n}$; in fact, $\omega$ having this value, any integer power $\omega^{k}$ is $= \cos\frac{2\pi k}{n} + i\sin\frac{2\pi k}{n}$, and we thence have $(\omega^{k})^{n} = \cos 2\pi k + i\sin 2\pi k = 1$, that is, $\omega^{k}$ is a root of the equation. The theory will be resumed further on.

\smallskip
\small\noindent\rule{0.4\textwidth}{0.4pt}\\
\textsuperscript{*} The square root of $\alpha + \beta i$ can be determined by the extraction of square roots of positive real numbers, without the trigonometrical tables.
\normalsize

By what precedes, we are led to the notion (a numerical) of the radical $a^{\frac{1}{n}}$ regarded as an $n$-valued function; any one of these being denoted by $\sqrt[n]{a}$, then the series of values is $\sqrt[n]{a}$, $\omega\sqrt[n]{a}$, \ldots, $\omega^{n-1}\sqrt[n]{a}$; or we may, if we please, use $\sqrt[n]{a}$ instead of $a^{\frac{1}{n}}$ as a symbol to denote the $n$-valued function.

As the coefficients of an algebraical equation may be numerical, all which follows in regard to algebraical equations is (with, it may be, some few modifications) applicable to numerical equations; and hence, concluding for the present this subject, it will be convenient to pass on to algebraical equations.

\subsection*{II. Algebraical Equations (\S\S 19 to 34).}

\subsection*{\S\,19.}

The equation is
\[
x^{n} - p_{1}x^{n-1} + \ldots \pm p_{n} = 0,
\]
and we here assume the existence of roots, viz.\ we assume that there are $n$ quantities $a, b, c, \ldots$ (in general all of them different, but which in particular cases may become equal in sets in any manner), such that
\[
x^{n} - p_{1}x^{n-1} + \ldots \pm p_{n} = 0;
\]
or looking at the question in a different point of view, and starting with the roots $a, b, c, \ldots$ as given, we express the product of the $n$ factors $x - a$, $x - b$, \ldots\ in the foregoing form, and thus arrive at an equation of the order $n$ having the $n$ roots $a, b, c, \ldots$ In either case we have
\[
p_{1} = \Sigma a, \quad p_{2} = \Sigma ab, \quad \ldots, \quad p_{n} = abc\ldots;
\]
i.e.\ regarding the coefficients $p_{1}, p_{2}, \ldots, p_{n}$ as given, then we assume the existence of roots $a, b, c, \ldots$ such that $p_{1} = \Sigma a$, \&c.; or, regarding the roots as given, then we write $p_{1}, p_{2}$, \&c., to denote the functions $\Sigma a$, $\Sigma ab$, \&c.

As already explained, the epithet algebraical is not used in opposition to numerical: an algebraical equation is merely an equation wherein the coefficients are not restricted to denote, or are not explicitly considered as denoting, numbers. That the abstraction is legitimate, appears by the simplest example; in saying that the equation $x^{2} - px + q = 0$ has a root $x = \frac{1}{2}(p + \sqrt{p^{2} - 4q})$, we mean that writing this value for $x$ the equation becomes an identity, $\frac{1}{2}(p + \sqrt{p^{2} - 4q})^{2} - p\cdot\frac{1}{2}(p + \sqrt{p^{2} - 4q}) + q = 0$; and the verification of this identity in nowise depends upon $p$ and $q$ meaning numbers. But if it be asked what there is beyond numerical equations included in the term algebraical equation, or, again, what is the full extent of the meaning attributed to the term: the latter question at any rate it would be very difficult to answer; as to the former one, it may be said that the coefficients may, for instance, be symbols of operation. As regards such equations, there is certainly no proof that every equation has a root, or that an equation of the $n$th order has $n$ roots; nor is it in any wise clear what the precise signification of the statement is. But it is found that the assumption of the existence of the $n$ roots can be made without contradictory results; conclusions derived from it, if they involve the roots, rest on the same ground as the original assumption; but the conclusion may be independent of the roots altogether, and in this case it is undoubtedly valid; the reasoning, although actually conducted by aid of the assumption (and, it may be, most easily and elegantly in this manner), is really independent of the assumption. In illustration, we observe that it is allowable to express a function of $p$ and $q$ as follows, that is, by means of a rational symmetrical function of $a$ and $b$: this can, as a fact, be expressed as a rational function of $a + b$ and $ab$; and if we prescribe that $a + b$ and $ab$ shall then be changed into $p$ and $q$ respectively, we have the required function of $p$, $q$. That is, we have $F(\alpha, \beta)$ as a representation of $f(p, q)$, obtained as if we had $p = a + b$, $q = ab$, but without in any wise assuming the existence of the $a$, $b$ of these equations.

\subsection*{\S\,20.}

Starting from the equation
\[
x^{n} - p_{1}x^{n-1} + \ldots = x - a \cdot x - b \cdot \&\text{c.},
\]
or the equivalent equations $p_{1} = \Sigma a$, \&c., we find
\[
a^{n} - p_{1}a^{n-1} + \ldots = 0, \quad b^{n} - p_{1}b^{n-1} + \ldots = 0;
\]
(it is as satisfying these equations that $a, b, \ldots$ are said to be the roots of $x^{n} - p_{1}x^{n-1} + \ldots = 0$) and conversely from the last-mentioned equations, assuming that $a, b, \ldots$ are all different, we deduce
\[
p_{1} = \Sigma a, \quad p_{2} = \Sigma ab, \quad \&\text{c.},
\]
and
\[
x^{n} - p_{1}x^{n-1} + \ldots = x - a \cdot x - b \cdot \&\text{c.}
\]

Observe that if, for instance, $a = b$, then the equations $a^{n} - p_{1}a^{n-1} + \ldots = 0$, $b^{n} - p_{1}b^{n-1} + \ldots = 0$ would reduce themselves to a single relation, which would not of itself express that $a$ was a double root,---that is, that $(x - a)^{2}$ was a factor of $x^{n} - p_{1}x^{n-1} + \&$c.; but by considering $b$ as the limit of $a + h$, $h$ indefinitely small, we obtain a second equation
\[
na^{n-1} - (n - 1)p_{1}a^{n-2} + \ldots = 0,
\]
which, with the first, expresses that $a$ is a double root; and then the whole system of equations leads as before to the equations $p_{1} = \Sigma a$, \&c.\ But the existence of a double root implies a certain relation between the coefficients; the general case is when the roots are all unequal.

We have then the theorem that every rational symmetrical function of the roots is a rational function of the coefficients. This is an easy consequence from the less general theorem, every rational and integral symmetrical function of the roots is a rational and integral function of the coefficients.

In particular, the sums of the powers $\Sigma a^{2}$, $\Sigma a^{3}$, \&c., are rational and integral functions of the coefficients.

The process originally employed for the expression of other functions $\Sigma a^{r}b^{s}$, \&c., in terms of the coefficients is to make them depend upon the sums of powers: for instance, $\Sigma a^{r}b^{s} = \Sigma a^{r} \cdot \Sigma a^{s} - \Sigma a^{r+s}$; but this is very objectionable; the true theory consists in showing that we have systems of equations
\begin{align*}
p_{1} &= \Sigma a, \\
p_{1}^{2} &= \Sigma a^{2} + 2\Sigma ab, \\
p_{2} &= \Sigma ab, \\
p_{3} &= \Sigma abc, \\
p_{1}p_{2} &= \Sigma a^{2}b + 3\Sigma abc, \\
p_{1}^{3} &= \Sigma a^{3} + 3\Sigma a^{2}b + 6\Sigma abc,
\end{align*}
where in each system there are precisely as many equations as there are root-functions on the right-hand side: e.g.\ 3 equations and 3 functions $\Sigma abc$, $\Sigma a^{2}b$, $\Sigma a^{3}$. Hence in each system the root-functions can be determined linearly in terms of the powers and products of the coefficients:
\begin{align*}
\Sigma ab &= p_{2}, \\
\Sigma abc &= p_{3}, \\
\Sigma a^{2}b &= p_{1}p_{2} - 3p_{3}, \\
\Sigma a^{3} &= p_{1}^{3} - 3p_{1}p_{2} + 3p_{3},
\end{align*}
and so on. The older process, if applied consistently, would derive the originally assumed value $\Sigma ab = p_{2}$, from the two equations $\Sigma a = p_{1}$, $\Sigma a^{2} = p_{1}^{2} - 2p_{2}$; i.e.\ we have $2\Sigma ab = \Sigma a \cdot \Sigma a - \Sigma a^{2} = p_{1}^{2} - (p_{1}^{2} - 2p_{2}) = 2p_{2}$.

\subsection*{\S\,21.}

It is convenient to mention here the theorem that, $x$ being determined as above by an equation of the order $n$, any rational and integral function whatever of $x$, or more generally any rational function which does not become infinite in virtue of the equation itself, can be expressed as a rational and integral function of $x$, of the order $n - 1$, the coefficients being rational functions of the coefficients of the equation. Thus the equation gives $x^{n}$ a function of the form in question; multiplying each side by $x$, and on the right-hand side writing for $x^{n}$ its foregoing value, we have $x^{n+1}$ a function of the form in question; and the like for any higher power of $x$, and therefore also for any rational and integral function of $x$. The proof in the case of a rational non-integral function is somewhat more complicated. The final result is of the form $\frac{\phi(x)}{\psi(x)} = I(x)$, or say $\phi(x) - \psi(x)I(x) = 0$, where $\phi$, $\psi$, $I$ are rational and integral functions; in other words, this equation, being true if only $f(x) = 0$, can only be so by reason that the left-hand side contains $f(x)$ as a factor, or we must have identically $\phi(x) - \psi(x)I(x) = M(x)f(x)$. And it is, moreover, clear that the equation $\frac{\phi(x)}{\psi(x)} = I(x)$, being satisfied if only $f(x) = 0$, must be satisfied by each root of the equation.

\subsection*{\S\,22.}

From the theorem that a rational symmetrical function of the roots is expressible in terms of the coefficients, it at once follows that it is possible to determine an equation (of an assignable order) having for its roots the several values of any given (unsymmetrical) function of the roots of the given equation. For example, in the case of a quartic equation, having the roots $(a, b, c, d)$, it is possible to find an equation having the roots $ab$, $ac$, $ad$, $bc$, $bd$, $cd$, being therefore a sextic equation: viz.\ in the product
\[
(y - ab)(y - ac)(y - ad)(y - bc)(y - bd)(y - cd),
\]
the coefficients of the several powers of $y$ will be symmetrical functions of $a, b, c, d$ and therefore rational and integral functions of the coefficients of the quartic equation; hence, supposing the product so expressed, and equating it to zero, we have the required sextic equation. In the same manner can be found the sextic equation having the roots $(a - b)^{2}$, $(a - c)^{2}$, $(a - d)^{2}$, $(b - c)^{2}$, $(b - d)^{2}$, $(c - d)^{2}$, which is the equation of differences previously referred to; and similarly we obtain the equation of differences for a given equation of any order. Again, the equation sought for may be that having for its roots the given rational functions $\phi(a)$, $\phi(b)$, \ldots\ of the several roots of the given equation. Any such rational function can (as was shown) be expressed as a rational and integral function of the order $n - 1$; and, retaining $x$ in place of any one of the roots, the problem is to find $y$ from the equations $x^{n} - p_{1}x^{n-1} + \ldots = 0$, and $y = M_{0}x^{n-1} + M_{1}x^{n-2} + \ldots$, or, what is the same thing, from these two equations to eliminate $x$. This is, in fact, Tschirnhausen's transformation (1683).

\subsection*{\S\,23.}

In connexion with what precedes, the question arises as to the number of values (obtained by permutations of the roots) of given unsymmetrical functions of the roots, or say of a given set of letters: for instance, with roots or letters $(a, b, c, d)$ as before, how many values are there of the function $ab + cd$, or better, how many functions are there of this form? The answer is 3, viz.\ $ab + cd$, $ac + bd$, $ad + bc$; or again we may ask whether, in the case of a given number of letters, there exist functions with a given number of values, 3-valued, 4-valued functions, \&c.

It is at once seen that for any given number of letters there exist 2-valued functions; the product of the differences of the letters is such a function; however the letters are interchanged, it alters only its sign; or say the two values are $\Delta$, $-\Delta$. And if $P$, $Q$ are symmetrical functions of the letters, then the general form of such a function is $P + Q\Delta$: this has only the two values $P + Q\Delta$, $P - Q\Delta$.

In the case of 4 letters there exist (as appears above) 3-valued functions: but in the case of 5 letters there does not exist any 3-valued or 4-valued function; and the only 5-valued functions are those which are symmetrical in regard to four of the letters, and can thus be expressed in terms of one letter and of symmetrical functions of all the letters. These last theorems present themselves in the demonstration of the non-existence of a solution of a quintic equation by radicals.

The theory is an extensive and important one, depending on the notions of substitutions and of groups\textsuperscript{*}.

\smallskip
\small\noindent\rule{0.4\textwidth}{0.4pt}\\
\textsuperscript{*} A \textit{substitution} is the operation by which we pass from the primitive arrangement of $n$ letters to any other arrangement of the same letters: for instance, the substitution $\binom{a}{b}$ means that $a$ is to be changed into $b$, $b$ into $c$, $c$ into $d$, $d$ into $a$. Substitutions may, of course, be represented by single letters $\alpha$, $\beta$, \ldots\ Two or more substitutions may be compounded together and give rise to a substitution; i.e., performing upon the primitive arrangement first the substitution $\beta$ and then upon the result the substitution $\alpha$, we have the substitution $\alpha\beta$. Substitutions are not commutative; thus, $\alpha\beta$ is not in general $= \beta\alpha$: but they are associative, $\alpha\beta \cdot \gamma = \alpha \cdot \beta\gamma$, so that $\alpha\beta\gamma$ has a determinate meaning. A substitution may be compounded any number of times with itself, and we thus have the powers $\alpha^{2}$, $\alpha^{3}$, \&c.\ Since the number of substitutions is limited, some power $\alpha^{\mu}$ must be $=1$: or, as this may be expressed, every substitution is a root of unity. A \textit{group} of substitutions is a set such that each two of them compounded together in either order gives a substitution belonging to the set; every group includes the substitution unity, so that we may in general speak of a group $1, \alpha, \beta, \ldots$ (the number of terms is the \textit{order} of the group). The whole system of the $1 \cdot 2 \cdot 3 \ldots n$ substitutions which can be performed upon the $n$ letters is obviously a group: the order of every other group which can be formed out of these substitutions is a submultiple of this number; but it is not conversely true that a group exists the order of which is any given submultiple of this number. In the case of a determinant the substitutions which give rise to the positive terms form a group the order of which is $= \frac{1}{2} \cdot 1 \cdot 2 \cdot 3 \ldots n$. For any function of the $n$ letters, the whole series of substitutions which leave the value of the functions unaltered form a group; and thence also the number of values of the function is $= 1 \cdot 2 \cdot 3 \ldots n$ divided by the order of the group.
\normalsize

\subsection*{\S\,24.}

Returning to equations, we have the very important theorem that, given the value of any unsymmetrical function of the roots, e.g.\ in the case of a quartic equation, the function $ab + cd$, it is in general possible to determine rationally the value of any similar function, such as $(a + b)^{2} + (c + d)^{2}$.

The \textit{a priori} ground of this theorem may be illustrated by means of a numerical equation. Suppose that the roots of a quartic equation are $1, 2, 3, 4$, then if it is given that $ab + cd = 14$, this in effect determines $a, b$ to be $1, 2$ and $c, d$ to be $3, 4$ (viz.\ $a = 1$, $b = 2$ or $a = 2$, $b = 1$, and $c = 3$, $d = 4$ or $c = 4$, $d = 3$) or else $a, b$ to be $3, 4$ and $c, d$ to be $1, 2$; and it therefore in effect determines $(a + b)^{2} + (c + d)^{2}$ to be $= 370$, and not any other value: that is, $(a + b)^{2} + (c + d)^{2}$, as having a single value, must be determinable rationally. And we can in the same way account for cases of failure as regards particular equations; thus, the roots being $1, 2, 3, 4$ as before, $a^{2}b = 2$ determines $a$ to be $= 1$ and $b$ to be $= 2$; but if the roots had been $1, 2, 4, 16$ then $a^{2}b = 16$ does not uniquely determine $a, b$ but only makes them to be $1, 16$ or $2, 4$ respectively.

As to the \textit{a posteriori} proof, assume, for instance,
\begin{align*}
t_{1} &= ab + cd, & y_{1} &= (a + b)^{2} + (c + d)^{2}, \\
t_{2} &= ac + bd, & y_{2} &= (a + c)^{2} + (b + d)^{2}, \\
t_{3} &= ad + bc, & y_{3} &= (a + d)^{2} + (b + c)^{2},
\end{align*}
then $y_{1} + y_{2} + y_{3}$, $t_{1}y_{1} + t_{2}y_{2} + t_{3}y_{3}$, $t_{1}^{2}y_{1} + t_{2}^{2}y_{2} + t_{3}^{2}y_{3}$, will be respectively symmetrical functions of the roots of the quartic, and therefore rational and integral functions of the coefficients; that is, they will be known.

Suppose for a moment that $t_{1}, t_{2}, t_{3}$ are all known; then the equations being linear in $y_{1}, y_{2}, y_{3}$, these can be expressed rationally in terms of the coefficients and of $t_{1}, t_{2}, t_{3}$; that is, $y_{1}, y_{2}, y_{3}$ will be known. But observe further that $y_{1}$ is obtained as a function of $t_{1}, t_{2}, t_{3}$ symmetrical as regards $t_{2}, t_{3}$; it can therefore be expressed as a rational function of $t_{1}$ and of $t_{2} + t_{3}$, $t_{2}t_{3}$, and thence as a rational function of $t_{1}$ and of $t_{1} + t_{2} + t_{3}$, $t_{1}t_{2} + t_{1}t_{3} + t_{2}t_{3}$, $t_{1}t_{2}t_{3}$; but these last are symmetrical functions of the roots, and as such they are expressible rationally in terms of the coefficients; that is, $y_{1}$ will be expressed as a rational function of $t_{1}$ and of the coefficients: or $t_{1}$ (alone, not $t_{2}$ or $t_{3}$) being known, $y_{1}$ will be rationally determined.

\subsection*{\S\,25.}

We now consider the question of the algebraical solution of equations, or, more accurately, that of the solution of equations by radicals.

In the case of a quadric equation $x^{2} - px + q = 0$, we can by the assistance of the sign $\surd(\quad)$ or $(\quad)^{\frac{1}{2}}$ find an expression for $x$ as a two-valued function of the coefficients $p, q$ such that, substituting this value in the equation, the equation is thereby identically satisfied; it has been found that this expression is
\[
x = \tfrac{1}{2}\{p \pm \sqrt{p^{2} - 4q}\},
\]
and the equation is on this account said to be algebraically solvable, or more accurately solvable by radicals. Or we may by writing $x = -\frac{1}{2}p + z$, reduce the equation to $z^{2} = \frac{1}{4}(p^{2} - 4q)$, viz.\ to an equation of the form $z^{2} = a$; and in virtue of its being thus reducible we say that the original equation is solvable by radicals. And the question for an equation of any higher order, say of the order $n$, is, can we by means of radicals, that is, by aid of the sign $\sqrt[m]{(\quad)}$ or $(\quad)^{\frac{1}{m}}$, using as many as we please of such signs and with any values of $m$, find an $n$-valued function (or any function) of the coefficients which substituted for $x$ in the equation shall satisfy it identically.

It will be observed that the coefficients $p, q, \ldots$ are not explicitly considered as numbers, but even if they do denote numbers, the question whether a numerical equation admits of solution by radicals is wholly unconnected with the before-mentioned theorem of the existence of the $n$ roots of such an equation. It does not even follow that in the case of a numerical equation solvable by radicals the algebraical solution gives the numerical solution, but this requires explanation. Consider first a numerical quadric equation with imaginary coefficients. In the formula $x = \frac{1}{2}(p \pm \sqrt{p^{2} - 4q})$, substituting for $p, q$ their given numerical values, we obtain for $x$ an expression of the form $x = \alpha + \beta i \pm \sqrt{\gamma + \delta i}$, where $\alpha, \beta, \gamma, \delta$ are real numbers. This expression substituted for $x$ in the quadric equation would satisfy it identically, and it is thus an algebraical solution; but there is no obvious \textit{a priori} reason why $\sqrt{\gamma + \delta i}$ should have a value $= c + di$, where $c$ and $d$ are real numbers calculable by the extraction of a root or roots of real numbers: however the case is (what there was no \textit{a priori} right to expect) that $\sqrt{\gamma + \delta i}$ has such a value calculable by means of the radical expressions $\sqrt[4]{\gamma^{2} + \delta^{2}} \pm \gamma$: and hence the algebraical solution of a numerical quadric equation does in every case give the numerical solution. The case of a numerical cubic equation will be considered presently.

\subsection*{\S\,26.}

A cubic equation can be solved by radicals. Taking for greater simplicity the cubic in the reduced form $x^{3} + qx - r = 0$, and assuming $x = a + b$, this will be a solution if only $3ab = q$ and $a^{3} + b^{3} = r$, equations which give $(a^{3} - b^{3})^{2} = r^{2} - \frac{4}{27}q^{3}$, a quadric equation solvable by radicals, and giving $a^{3} - b^{3} = \sqrt{r^{2} - \frac{4}{27}q^{3}}$, a two-valued function of the coefficients: combining this with $a^{3} + b^{3} = r$, we have $a^{3} = \frac{1}{2}(r + \sqrt{r^{2} - \frac{4}{27}q^{3}})$, a two-valued function: we then have $a$ by means of a cube root, viz.
\[
a = \sqrt[3]{\tfrac{1}{2}(r + \sqrt{r^{2} - \tfrac{4}{27}q^{3}})},
\]
a six-valued function of the coefficients; but then, writing $q = 3ab$, we have, as may be shown, $a + b$ a three-valued function of the coefficients; and $x = a + b$ is the required solution by radicals. It would have been wrong to complete the solution by writing
\[
b = \sqrt[3]{\tfrac{1}{2}(r - \sqrt{r^{2} - \tfrac{4}{27}q^{3}})},
\]
for then $a + b$ would have been given as a 9-valued function having only 3 of its values roots, and the other 6 values being irrelevant. Observe that in this last process we make no use of the equation $3ab = q$, in its original form, but use only the derived equation $27a^{3}b^{3} = q^{3}$, implied in, but not implying, the original form.

An interesting variation of the solution is to write $x = ab(a + b)$, giving $a^{3}b^{3}(a^{3} + b^{3}) = \frac{q^{3}}{27}$ and $3a^{3}b^{3} = q^{3}$, or say $a^{3} + b^{3} = \frac{3r}{q^{3}}$, $a^{3}b^{3} = \frac{q^{3}}{27}$; and consequently
\[
a^{3} = \tfrac{1}{2}\left(\frac{3r}{q^{3}} + \sqrt{\frac{9r^{2}}{q^{6}} - \frac{4q^{3}}{27}}\right), \quad b^{3} = \tfrac{1}{2}\left(\frac{3r}{q^{3}} - \sqrt{\frac{9r^{2}}{q^{6}} - \frac{4q^{3}}{27}}\right),
\]
i.e., here $a^{3}$, $b^{3}$ are each of them a two-valued function, but as the only effect of altering the sign of the quadric radical is to interchange $a^{3}$, $b^{3}$, they may be regarded as each of them one-valued; $a$ and $b$ are each of them 3-valued (for observe that here only $a^{3}b^{3}$, not $ab$, is given); and $ab(a + b)$ thus is in appearance a 9-valued function, but it can easily be shown that it is (as it ought to be) only 3-valued.

In the case of a numerical cubic, even when the coefficients are real, substituting their values in the expression
\[
x = \sqrt[3]{\tfrac{1}{2}(r + \sqrt{r^{2} - \tfrac{4}{27}q^{3}})} + \sqrt[3]{\tfrac{1}{2}(r - \sqrt{r^{2} - \tfrac{4}{27}q^{3}})},
\]
this may depend on an expression of the form $\sqrt[3]{\gamma + \delta i}$, where $\gamma$ and $\delta$ are real numbers (it will do so if $r^{2} - \frac{4}{27}q^{3}$ is a negative number), and then we cannot by the extraction of any root or roots of real positive numbers reduce $\sqrt[3]{\gamma + \delta i}$ to the form $c + di$, $c$ and $d$ real numbers; hence here the algebraical solution does not give the numerical solution, and we have here the so-called ``irreducible case'' of a cubic equation. By what precedes, there is nothing in this that might not have been expected: the algebraical solution makes the solution depend on the extraction of the cube root of a negative number, and there was no reason for expecting this to be a real number. It is well known that the case in question is that wherein the three roots of the numerical cubic equation are all real; if the roots are two imaginary, one real, then contrariwise the quantity under the cube root is real; and the algebraical solution gives the numerical one.

The irreducible case is solvable by a trigonometrical formula, but this is not a solution by radicals: it consists, in effect, in reducing the given numerical cubic (not to a cubic of the form $z^{3} = a$, solvable by the extraction of a cube root, but) to a cubic of the form $4z^{3} - 3z = a$, corresponding to the equation $4\cos^{3}\theta - 3\cos\theta = \cos 3\theta$ which serves to determine $\cos\theta$ when $\cos 3\theta$ is known. The theory is applicable to an algebraical cubic equation: say that such an equation, if it can be reduced to the form $4z^{3} - 3z = a$, is solvable by ``trisection'': then the general cubic equation is solvable by trisection.

\subsection*{\S\,27.}

A quartic equation is solvable by radicals: and it is to be remarked that the existence of such a solution depends on the existence of 3-valued functions such as $ab + cd$ of the four roots $(a, b, c, d)$: by what precedes, $ab + cd$ is the root of a cubic equation, which equation is solvable by radicals: hence $ab + cd$ can be found by radicals; and since $abcd$ is a given function, $ab$ and $cd$ can then be found by radicals. But by what precedes, if $ab$ be known then any similar function, say $a + b$, is obtainable rationally; and then from the values of $a + b$ and $ab$ we may by radicals obtain the value of $a$ or $b$, that is, an expression for the root of the given quartic equation: the expression ultimately obtained is 4-valued, corresponding to the different values of the several radicals which enter therein, and we have thus the expression by radicals of each of the four roots of the quartic equation. But when the quartic is numerical the same thing happens as in the cubic, and the algebraical solution does not in every case give the numerical one.

It will be understood, from the foregoing explanation as to the quartic, how in the next following case, that of the quintic, the question of the solvability by radicals depends on the existence or non-existence of 5-valued functions of the five roots $(a, b, c, d, e)$; the fundamental theorem is the one already stated, a rational function of five letters, if it has less than 5, cannot have more than 2 values, that is, there are no 3-valued or 4-valued functions of 5 letters; and by reasoning depending in part upon this theorem, Abel (1824) showed that a general quintic equation is not solvable by radicals; and \textit{a fortiori} the general equation of any order higher than 5 is not solvable by radicals.

\subsection*{\S\,28.}

The general theory of the solvability of an equation by radicals depends fundamentally on Vandermonde's remark (1770) that, supposing an equation is solvable by radicals, and that we have therefore an algebraical expression of $x$ in terms of the coefficients, then substituting for the coefficients their values in terms of the roots, the resulting expression must reduce itself to any one at pleasure of the roots $a, b, c, \ldots$; thus in the case of the quadric equation, in the expression $x = \frac{1}{2}(p + \sqrt{p^{2} - 4q})$, substituting for $p$ and $q$ their values, and observing that $(a + b)^{2} - 4ab = (a - b)^{2}$, this becomes $x = \frac{1}{2}\{a + b + \sqrt{(a - b)^{2}}\}$, the value being $a$ or $b$ according as the radical is taken to be $+(a - b)$ or $-(a - b)$.

So in the cubic equation $x^{3} - px^{2} + qx - r = 0$, if the roots are $a, b, c$, and if $\omega$ is used to denote an imaginary cube root of unity, $\omega^{2} + \omega + 1 = 0$, then writing for shortness $p = a + b + c$, $L = a + \omega b + \omega^{2}c$, $M = a + \omega^{2}b + \omega c$, it is at once seen that $LM$, $L^{3} + M^{3}$, and therefore also $(L^{3} - M^{3})^{2}$ are symmetrical functions of the roots, and consequently rational functions of the coefficients: hence
\[
\tfrac{1}{2}\{L^{3} + M^{3} + \sqrt{(L^{3} - M^{3})^{2}}\}
\]
is a rational function of the coefficients, which when these are replaced by their values as functions of the roots becomes, according to the sign given to the quadric radical, $= L^{3}$ or $M^{3}$: taking it $= L^{3}$, the cube root of the expression has the three values $L$, $\omega L$, $\omega^{2}L$; and $LM$ divided by the same cube root has therefore the values $M$, $\omega^{2}M$, $\omega M$; whence finally the expression
\[
\tfrac{1}{3}\{p + LM^{2} \div \sqrt[3]{\tfrac{1}{2}(L^{3} + M^{3} + \sqrt{(L^{3} - M^{3})^{2}})}\}
\]
has the three values
\[
\tfrac{1}{3}(p + L + M), \quad \tfrac{1}{3}(p + \omega L + \omega^{2}M), \quad \tfrac{1}{3}(p + \omega^{2}L + \omega M);
\]
that is, these are $= a$, $b$, $c$ respectively. If the value $M^{3}$ had been taken instead of $L^{3}$, then the expression would have had the same three values $a, b, c$. Comparing the solution given for the cubic $x^{3} + qx - r = 0$, it will readily be seen that the two solutions are identical, and that the function $r^{2} - \frac{4}{27}q^{3}$ under the radical sign must (by aid of the relation $p = 0$ which subsists in this case) reduce itself to $(L^{3} - M^{3})^{2}$; it is only by each radical being equal to a rational function of the roots that the final expression can become equal to the roots $a, b, c$ respectively.

\subsection*{\S\,29.}

The formulae for the cubic were obtained by Lagrange (1770--71) from a different point of view. Upon examining and comparing the principal known methods for the solution of algebraical equations, he found that they all ultimately depended upon finding a ``resolvent'' equation of which the root is $a + \omega b + \omega^{2}c + \omega^{3}d + \ldots$, $\omega$ being an imaginary root of unity, of the same order as the equation: e.g.\ for the cubic the root is $a + \omega b + \omega^{2}c$, $\omega$ an imaginary cube root of unity. Evidently the method gives for $L^{3}$ a quadric equation, which is the ``resolvent'' equation in this particular case.

For a quartic the formulae present themselves in a somewhat different form, by reason that 4 is not a prime number. Attempting to apply it to a quintic, we seek for the equation of which the root is $(a + \omega b + \omega^{2}c + \omega^{3}d + \omega^{4}e)$, $\omega$ an imaginary fifth root of unity, or rather the fifth power thereof $(a + \omega b + \omega^{2}c + \omega^{3}d + \omega^{4}e)^{4}$; this is a 24-valued function, but if we consider the four values corresponding to the roots of unity $\omega$, $\omega^{2}$, $\omega^{3}$, $\omega^{4}$, viz.\ the values
\begin{align*}
&(a + \omega b + \omega^{2}c + \omega^{3}d + \omega^{4}e)^{4}, \\
&(a + \omega^{2}b + \omega^{4}c + \omega d + \omega^{3}e)^{4}, \\
&(a + \omega^{3}b + \omega c + \omega^{4}d + \omega^{2}e)^{4}, \\
&(a + \omega^{4}b + \omega^{3}c + \omega^{2}d + \omega e)^{4},
\end{align*}
any symmetrical function of these, for instance their sum, is a six-valued function of the roots, and may therefore be determined by means of a sextic equation, the coefficients whereof are rational functions of the coefficients of the original quintic equation; the conclusion being that the solution of an equation of the fifth order is made to depend upon that of an equation of the sixth order. This is, of course, useless for the solution of the quintic equation, which, as already mentioned, does not admit of solution by radicals; but the equation of the sixth order, Lagrange's resolvent sextic, is very important, and is intimately connected with all the later investigations in the theory.

\subsection*{\S\,30.}

It is to be remarked, in regard to the question of solvability by radicals, that not only the coefficients are taken to be arbitrary, but it is assumed that they are represented each by a single letter, or say rather that they are not so expressed in terms of other arbitrary quantities as to make a solution possible. If the coefficients are not all arbitrary, for instance, if some of them are zero, a sextic equation might be of the form $ax^{6} + bx^{4} + cx^{2} + d = 0$, and so be solvable as a cubic; or if the coefficients of the sextic are given functions of the six arbitrary quantities $a, b, c, d, e, f$, such that the sextic is really of the form
\[
(x^{2} + ax + b)(x^{4} + cx^{3} + dx^{2} + ex + f) = 0,
\]
then it breaks up into the equations $x^{2} + ax + b = 0$, $x^{4} + cx^{3} + dx^{2} + ex + f = 0$, and is consequently solvable by radicals; so also if the form is
\[
(x - a)(x - b)(x - c)(x - d)(x - e)(x - f) = 0,
\]
then the equation is solvable by radicals,---in this extreme case rationally. Such cases of solvability are self-evident; but they are enough to show that the general theorem of the non-solvability by radicals of an equation of the fifth or any higher order does not in any wise exclude for such orders the existence of particular equations solvable by radicals, and there are, in fact, extensive classes of equations which are thus solvable; the binomial equations $x^{n} - 1 = 0$ present an instance.

It has already been shown how the several roots of the equation $x^{n} - 1 = 0$ can be expressed in the form $\cos\frac{2\pi}{n} + i\sin\frac{2\pi}{n}$, but the question is now that of the algebraical solution (or solution by radicals) of this equation. There is always a root $= 1$; if $\omega$ be any other root, then obviously $\omega, \omega^{2}, \ldots, \omega^{n-1}$ are all of them roots: $x^{n} - 1$ contains the factor $x - 1$, and it thus appears that $\omega, \omega^{2}, \ldots, \omega^{n-1}$ are the $n - 1$ roots of the equation
\[
x^{n-1} + x^{n-2} + \ldots + x + 1 = 0;
\]
we have, of course,
\[
\omega^{n-1} + \omega^{n-2} + \ldots + \omega + 1 = 0.
\]

It is proper to distinguish the cases $n$ prime and $n$ composite; and in the latter case there is a distinction of $n$ simple or multiple according as the prime factors are.

By way of illustration, suppose successively $n = 15$ and $n = 9$: in the former case, if $\alpha$ be an imaginary root of $x^{3} - 1 = 0$ (or root of $x^{2} + x + 1 = 0$), and $\beta$ an imaginary root of $x^{5} - 1 = 0$ (or root of $x^{4} + x^{3} + x^{2} + x + 1 = 0$), then $\omega$ may be taken $= \alpha\beta$; the successive powers thereof, $\alpha\beta$, $\alpha^{2}\beta^{2}$, $\beta^{3}$, $\alpha\beta^{4}$, $\alpha^{2}\beta^{5}$, $\alpha\beta^{6}$, $\alpha^{2}\beta^{7}$, $\beta^{8}$, $\alpha$, $\alpha^{2}\beta$, $\beta^{2}$, are the roots of $x^{14} + x^{13} + \ldots + x + 1 = 0$; the solution thus depends on the solution of the equations $x^{3} - 1 = 0$ and $x^{5} - 1 = 0$. In the latter case, if $\alpha$ be an imaginary root of $x^{3} - 1 = 0$ (or root of $x^{2} + x + 1 = 0$), then the equation $x^{9} - 1 = 0$ gives $x^{3} = 1$, $\alpha$, or $\alpha^{2}$; $x^{3} = 1$ gives $x = 1$, $\alpha$, or $\alpha^{2}$; and the solution thus depends on the solution of the equations $x^{3} - 1 = 0$, $x^{3} - \alpha = 0$, $x^{3} - \alpha^{2} = 0$. The first equation has the roots $1$, $\alpha$, $\alpha^{2}$; if $\beta$ be a root of either of the others, say if $\beta^{3} = \alpha$, then assuming $\omega = \beta$, the successive powers are $\beta$, $\beta^{2}$, $\alpha$, $\alpha\beta$, $\alpha\beta^{2}$, $\alpha^{2}$, $\alpha^{2}\beta$, $\alpha^{2}\beta^{2}$, which are the roots of the equation $x^{8} + x^{7} + \ldots + x + 1 = 0$.

It thus appears that the only case which need be considered is that of $n$ a prime number, and writing (as is more usual) $r$ in place of $\omega$, we have $r$, $r^{2}$, $r^{3}$, \ldots, $r^{n-1}$ as the $(n - 1)$ roots of the reduced equation
\[
x^{n-1} + x^{n-2} + \ldots + x + 1 = 0;
\]
then not only $r^{n} - 1 = 0$, but also $r^{n-1} + r^{n-2} + \ldots + r + 1 = 0$.

\subsection*{\S\,31.}

The process of solution due to Gauss (1801) depends essentially on the arrangement of the roots in a certain order, viz.\ not as above, with the indices of $r$ in arithmetical progression, but with their indices in geometrical progression; the prime number $n$ has a certain number of prime roots $g$, which are such that $g^{n-1}$ is the lowest power of $g$ which is $= 1$ to the modulus $n$: or, what is the same thing, that the series of powers $1, g, g^{2}, \ldots, g^{n-2}$, each divided by $n$, leave (in a different order) the remainders $1, 2, 3, \ldots, n - 1$; hence giving to $r$ in succession the indices $1, g, g^{2}, \ldots, g^{n-2}$, we have, in a different order, the whole series of roots $r, r^{g}, r^{g^{2}}, \ldots, r^{g^{n-2}}$.

In the most simple case, $n = 5$, the equation to be solved is $x^{4} + x^{3} + x^{2} + x + 1 = 0$; here 2 is a prime root of 5, and the order of the roots is $r$, $r^{2}$, $r^{4}$, $r^{3}$. The Gaussian process consists in forming an equation for determining the periods $P_{1}$, $P_{2} = r + r^{4}$ and $r^{2} + r^{3}$ respectively, these being such that the symmetrical functions $P_{1} + P_{2}$, $P_{1}P_{2}$ are rationally determinable: in fact,
\[
P_{1} + P_{2} = -1, \quad P_{1}P_{2} = (r + r^{4})(r^{2} + r^{3}) = r^{3} + r^{4} + r^{6} + r^{7} = r^{3} + r^{4} + r + r^{2} = -1.
\]
$P_{1}$, $P_{2}$ are thus the roots of $u^{2} + u - 1 = 0$; and taking them to be known, they are themselves broken up into subperiods, in the present case single terms, $r$ and $r^{4}$ for $P_{1}$, $r^{2}$ and $r^{3}$ for $P_{2}$; the symmetrical functions of these are then rationally determined in terms of $P_{1}$ and $P_{2}$; thus $r + r^{4} = P_{1}$, $r \cdot r^{4} = 1$, or $r$, $r^{4}$ are the roots of $x^{2} - P_{1}x + 1 = 0$. The mode of division is more clearly seen for a larger value of $n$; thus, for $n = 7$ a prime root is $= 3$, and the arrangement of the roots is $r$, $r^{3}$, $r^{2}$, $r^{6}$, $r^{4}$, $r^{5}$. We may form either 3 periods each of 2 terms, $P_{1}, P_{2}, P_{3} = r + r^{6}$, $r^{3} + r^{4}$, $r^{2} + r^{5}$ respectively; or else 2 periods each of 3 terms, $P_{1}, P_{2} = r + r^{3} + r^{2}$, $r^{6} + r^{4} + r^{5}$ respectively; in each case the symmetrical functions of the periods are rationally determinable; thus in the case of the two periods $P_{1} + P_{2} = -1$, $P_{1}P_{2} = 3 + r + r^{3} + r^{2} + r^{6} + r^{4} + r^{5} = 2$; and, the periods being known, the symmetrical functions of the several terms of each period are rationally determined in terms of the periods, thus
\[
r + r^{3} + r^{2} = P_{1}, \quad r \cdot r^{3} + r \cdot r^{2} + r^{3} \cdot r^{2} = P_{2}, \quad r \cdot r^{3} \cdot r^{2} = 1.
\]

The theory was further developed by Lagrange (1808), who, applying his general process to the equation in question, $x^{n-1} + x^{n-2} + \ldots + x + 1 = 0$, the roots $a, b, c, \ldots$ being the several powers of $r$, the indices in geometrical progression as above, showed that the function $(a + \omega b + \omega^{2}c + \ldots)^{n-1}$ was in this case a given function of $\omega$ with integer coefficients. Reverting to the before-mentioned particular equation $x^{4} + x^{3} + x^{2} + x + 1 = 0$, it is very interesting to compare the process of solution with that for the solution of the general quartic the roots whereof are $a, b, c, d$.

Take $\omega$, a root of the equation $\omega^{4} - 1 = 0$ (whence $\omega$ is $= 1$, $-1$, $i$, or $-i$, at pleasure), and consider the expression
\[
(a + \omega b + \omega^{2}c + \omega^{3}d)^{4}.
\]
The developed value of this is
\begin{multline*}
a^{4} + b^{4} + c^{4} + d^{4} + 6(a^{2}b^{2} + c^{2}d^{2}) + 12(a^{2}bd + b^{2}ca + c^{2}db + d^{2}ac) \\
+ \omega\{4(a^{3}b + b^{3}c + c^{3}d + d^{3}a) + 12(a^{2}cd + b^{2}da + c^{2}ab + d^{2}bc)\} \\
+ \omega^{2}\{6(a^{2}c^{2} + b^{2}d^{2}) + 4(a^{3}c + b^{3}d + c^{3}a + d^{3}b) + 24abcd\} \\
+ \omega^{3}\{4(a^{3}d + b^{3}a + c^{3}b + d^{3}c) + 12(a^{2}bc + b^{2}cd + c^{2}da + d^{2}ab)\};
\end{multline*}
that is, this is a 6-valued function of $a, b, c, d$, the root of a sextic (which is, in fact, solvable by radicals; but this is not here material).

If, however, $a, b, c, d$ denote the roots $r, r^{2}, r^{4}, r^{3}$ of the special equation, then the expression becomes
\begin{multline*}
r^{4} + r^{8} + r + r^{9} + 6(1 + 1) + 12(r^{6} + r^{4} + r^{7} + r^{3}) \\
+ \omega\{4(1 + 1 + 1 + 1) + 12(r^{6} + r^{4} + r + r^{9})\} \\
+ \omega^{2}\{6(r + r^{8} + r^{4} + r^{3}) + 4(r^{6} + r^{4} + r^{7} + r^{3})\} \\
+ \omega^{3}\{4(r + r^{8} + r^{4} + r^{9}) + 12(r^{6} + r + r^{3} + r^{4})\},
\end{multline*}
viz.\ this is $= -1 + 4\omega + 14\omega^{2} - 16\omega^{3}$, a completely determined value. That is, we have
\[
(r + \omega r^{2} + \omega^{2}r^{4} + \omega^{3}r^{3})^{4} = -1 + 4\omega + 14\omega^{2} - 16\omega^{3},
\]
which result contains the solution of the equation. If $\omega = 1$, we have $(r + r^{2} + r^{4} + r^{3})^{4} = 1$, which is right; if $\omega = -1$, then $(r + r^{4} - r^{2} - r^{3})^{4} = 25$; if $\omega = i$, then we have $\{r - r^{4} + i(r^{2} - r^{3})\}^{4} = -15 + 20i$; and if $\omega = -i$, then $\{r - r^{4} - i(r^{2} - r^{3})\}^{4} = -15 - 20i$; the solution may be completed without difficulty.

The result is perfectly general, thus: $n$ being a prime number, $r$ a root of the equation $x^{n-1} + x^{n-2} + \ldots + x + 1 = 0$, $\omega$ a root of $\omega^{n-1} - 1 = 0$, and $g$ a prime root of $g^{n-1} \equiv 1 \pmod{n}$, then
\[
(r + \omega r^{g} + \ldots + \omega^{n-2}r^{g^{n-2}})^{n-1}
\]
is a given function $M_{0} + M_{1}\omega + \ldots + M_{n-2}\omega^{n-2}$ with integer coefficients, and by the extraction of $(n - 1)$th and similar roots we ultimately obtain $r$ in terms of $\omega$, which is taken to be known: the equation $x^{n} - 1 = 0$, $n$ a prime number, is thus solvable by radicals. In particular, if $n - 1$ be a power of 2, the solution (by either process) requires the extraction of square roots only: and it was thus that Gauss discovered that it was possible to construct geometrically the regular polygons of 17 sides and 257 sides respectively. Some interesting developments in regard to the theory were obtained by Jacobi (1837); see the memoir ``Ueber die Kreistheilung, u.s.w.,'' \textit{Crelle}, t.~xxx.\ (1846).

\subsection*{\S\,32.}

The equation $x^{n-1} + \ldots + x + 1 = 0$ has been considered for its own sake, but it also serves as a specimen of a class of equations solvable by radicals, considered by Abel (1828), and since called \textit{Abelian equations}, viz., for the Abelian equation of the order $n$, if $x$ be any root, the roots are $x$, $\theta x$, $\theta^{2}x$, \ldots, $\theta^{n-1}x$ ($\theta x$ being a rational function of $x$, and $\theta^{n}x = x$); the theory is, in fact, very analogous to that of the above particular case. A more general theorem obtained by Abel is as follows:---If the roots of an equation of any order are connected together in such wise that all the roots can be expressed rationally in terms of any one of them, say $x$; if, moreover, $\theta x$, $\theta^{2}x$ being any two of the roots, we have $\theta\theta^{2}x = \theta^{2}\theta x$, the equation will be solvable algebraically. It is proper to refer also to Abel's definition of an irreducible equation:---an equation $\phi x = 0$, the coefficients of which are rational functions of a certain number of known quantities $a, b, c, \ldots$, is called irreducible when it is impossible to express its roots by an equation of an inferior degree, the coefficients of which are also rational functions of $a, b, c, \ldots$ (or, what is the same thing, when $\phi x$ does not break up into factors which are rational functions of $a, b, c, \ldots$). Abel applied his theory to the equations which present themselves in the division of the elliptic functions, but not to the modular equations.

\subsection*{\S\,33.}

But the theory of the algebraical solution of equations in its most complete form was established by Galois (born October 1811, killed in a duel May 1832; see his collected works, \textit{Liouville}, t.~xi., 1846). The definition of an irreducible equation resembles Abel's,---an equation is reducible when it admits of a rational divisor, irreducible in the contrary case; only the word rational is used in this extended sense that, in connexion with the coefficients of the given equation, or with the irrational quantities (if any) whereof these are composed, he considers any number of other irrational quantities called ``adjoint radicals,'' and he terms rational any rational function of the coefficients (or the irrationals whereof they are composed) and of these adjoint radicals; the epithet irreducible is thus taken either absolutely or in a relative sense, according to the system of adjoint radicals which are taken into account. For instance, the equation $x^{4} + x^{3} + x^{2} + x + 1 = 0$; the left-hand side has here no rational divisor, and the equation is irreducible; but this function is $= (x^{2} + \frac{1}{2}x + 1)^{2} - \frac{5}{4}x^{2}$, and it has thus the irrational divisors $x^{2} + \frac{1}{2}(1 + \sqrt{5})x + 1$, $x^{2} + \frac{1}{2}(1 - \sqrt{5})x + 1$ and these, if we adjoin the radical $\sqrt{5}$, are rational, and the equation is no longer irreducible.

In the case of a given equation, assumed to be irreducible, the problem to solve the equation is, in fact, that of finding radicals by the adjunction of which the equation becomes reducible; for instance, the general quadric equation $x^{2} + px + q = 0$ is irreducible, but it becomes reducible, breaking up into rational linear factors, when we adjoin the radical $\sqrt{\frac{1}{4}p^{2} - q}$.

The fundamental theorem is the Proposition I.\ of the ``M\'{e}moire sur les conditions de resolubilit\'{e} des Equations par radicaux''; viz.\ given an equation of which $a, b, c, \ldots$ are the roots, there is always a group of permutations of the letters $a, b, c, \ldots$ possessed of the following properties:
\begin{enumerate}
\item Every function of the roots invariable by the substitutions of the group is rationally known.
\item Reciprocally, every rationally determinable function of the roots is invariable by the substitutions of the group.
\end{enumerate}

Here by an invariable function is meant not only a function of which the form is invariable by the substitutions of the group, but further, one of which the value is invariable by these substitutions: for instance, if the equation be $\phi x = 0$, then $\phi x$ is a function of the roots invariable by any substitution whatever. And in saying that a function is rationally known, it is meant that its value is expressible rationally in terms of the coefficients and of the adjoint quantities.

For instance, in the case of a general equation, the group is simply the system of the $1 \cdot 2 \cdot 3 \ldots n$ permutations of all the roots, since, in this case, the only rationally determinable functions are the symmetric functions of the roots.

In the case of the equation $x^{n-1} + \ldots + x + 1 = 0$, $n$ a prime number, $a, b, c, \ldots, k = r, r^{g}, r^{g^{2}}, \ldots, r^{g^{n-2}}$, where $g$ is a prime root of $n$, then the group is the cyclical group $abc\ldots k$, $bc\ldots ka$, \ldots, $kab\ldots j$, that is, in this particular case the number of the permutations of the group is equal to the order of the equation.

This notion of the group of the original equation, or of the group of the equation as varied by the adjunction of a series of radicals, seems to be the fundamental one in Galois's theory. But the problem of solution by radicals, instead of being the sole object of the theory, appears as the first link of a long chain of questions relating to the transformation and classification of irrationals.

Returning to the question of solution by radicals, it will be readily understood that by the adjunction of a radical the group may be diminished; for instance, in the case of the general cubic, where the group is that of the six permutations, by the adjunction of the square root which enters into the solution, the group is reduced to $abc$, $bca$, $cab$; that is, it becomes possible to express rationally, in terms of the coefficients and of the adjoint square root, any function such as $a^{2}b + b^{2}c + c^{2}a$ which is not altered by the cyclical substitution $a$ into $b$, $b$ into $c$, $c$ into $a$. And hence, to determine whether an equation of a given form is solvable by radicals, the course of investigation is to inquire whether, by the successive adjunction of radicals, it is possible to reduce the original group of the equation so as to make it ultimately consist of a single permutation.

The condition in order that an equation of a given prime order $n$ may be solvable by radicals was in this way obtained---in the first instance in the form, scarcely intelligible without further explanation, that every function of the roots $x_{1}, x_{2}, \ldots, x_{n}$, invariable by the substitutions $x_{s} \mapsto ax_{s} + b$, must be rationally known; and then in the equivalent form that the resolvent equation of the order $1 \cdot 2 \cdot \ldots \cdot (n - 2)$ must have a rational root. In particular, the condition in order that a quintic equation may be solvable is that Lagrange's resolvent of the order 6 may have a rational factor, a result obtained from a direct investigation in a valuable memoir by E.\ Luther, \textit{Crelle}, t.~xxxiv.\ (1847).

Among other results demonstrated or announced by Galois may be mentioned those relating to the modular equations in the theory of elliptic functions; for the transformations of the orders 5, 7, 11, the modular equations of the orders 6, 8, 12 are depressible to the orders 5, 7, 11 respectively; but for the transformation, $n$ a prime number greater than 11, the depression is impossible.

The general theory of Galois in regard to the solution of equations was completed, and some of the demonstrations supplied, by Betti (1852). See also Serret's \textit{Cours d'Alg\`{e}bre sup\'{e}rieure}, 2nd ed.\ 1854; 4th ed.\ 1877--78.

\subsection*{\S\,34.}

Returning to quintic equations, Jerrard (1835) established the theorem that the general quintic equation is, by the extraction of only square and cubic roots, reducible to the form $x^{5} + ax + b = 0$, or what is the same thing, to $x^{5} + x + b = 0$. The actual reduction by means of Tschirnhausen's theorem was effected by Hermite in connexion with his elliptic-function solution of the quintic equation (1858) in a very elegant manner. It was shown by Cockle and Harley (1858--59) in connexion with the Jerrardian form, and by Cayley (1861), that Lagrange's resolvent equation of the sixth order can be replaced by a more simple sextic equation occupying a like place in the theory.

The theory of the modular equations, more particularly for the case $n = 5$, has been studied by Hermite, Kronecker, and Brioschi. In the case $n = 5$, the modular equation of the order 6 depends, as already mentioned, on an equation of the order 5 and conversely the general quintic equation may be made to depend upon this modular equation of the order 6: that is, assuming the solution of this modular equation, we can solve (not by radicals) the general quintic equation; this is Hermite's solution of the general quintic equation by elliptic functions (1858); it is analogous to the before-mentioned trigonometrical solution of the cubic equation. The theory is reproduced and developed in Brioschi's memoir, ``Ueber die Aufl\"{o}sung der Gleichungen vom f\"{u}nften Grade,'' \textit{Math.\ Annalen}, t.~xiii.\ (1877--78).

The great modern work, reproducing the theories of Galois, and exhibiting the theory of algebraic equations as a whole, is Jordan's \textit{Trait\'{e} des Substitutions et des Equations Alg\'{e}briques}, Paris, 1870. The work is divided into four books---book I., preliminary, relating to the theory of congruences; book II.\ is in two chapters, the first relating to substitutions in general, the second to substitutions defined analytically, and chiefly to linear substitutions; book III.\ has four chapters, the first discussing the principles of the general theory, the other three containing applications to algebra, geometry, and the theory of transcendents; book IV., divided into seven chapters, contains lastly, a determination of the general types of equations solvable by radicals, and a complete system of classification of these types. A glance through the index will show the vast extent which the theory has assumed, and the form of general conclusions arrived at; thus, in book III., the algebraical applications comprise Abelian equations, equations of Galois; the geometrical ones comprise Hesse's equation, Clebsch's equations, lines on a quartic surface having a nodal line, singular points of Kummer's surface, lines on a cubic surface, problems of contact; the applications to the theory of transcendents comprise circular functions, elliptic functions (including division and the modular equation), hyperelliptic functions, solution of equations by transcendents. And on this last subject, solution of equations by transcendents, we may quote the result, ``the solution of the general equation of an order superior to five cannot be made to depend upon that of the equations for the division of the circular or elliptic functions''; and again (but with a reference to a possible case of exception), ``the general equation cannot be solved by aid of the equations which give the division of the hyperelliptic functions into an odd number of parts.''

\newpage
% ============================================================
\section*{\S\,787. FUNCTION.}
\addcontentsline{toc}{section}{787. Function}
% ============================================================

\noindent[From the \textit{Encyclop{\ae}dia Britannica}, Ninth Edition, vol.~ix (1879), pp.~818--824.]

\smallskip

Functionality, in Analysis, is dependence on a variable or variables; in the case of a single variable $u$, it is the same thing to say that $v$ depends upon $u$, or to say that $v$ is a function of $u$, only in the latter form of expression the mode of dependence is embodied in the term ``function.'' We have given or known functions such as $u^{2}$ or $\sin u$, and the general notation of the form $\phi u$, where the letter $\phi$ is used as a functional symbol to denote a function of $u$, known or unknown as the case may be: in each case $u$ is the independent variable or argument of the function, but it is to be observed that, if $v$ be a function of $u$, then $v$ like $u$ is a variable, the values of $v$ regarded as known serve to determine those of $u$: that is, we may conversely regard $u$ as a function of $v$. In the case of two or more independent variables, say when $w$ depends on or is a function of $u$, $v$, \&c., or $w = \phi(u, v, \ldots)$, then $u, v, \ldots$ are the independent variables or arguments of the function; frequently when one of these variables, say $u$, is principally or alone attended to, it is regarded as the independent variable or argument of the function, and the other variables $v$, \&c., are regarded as parameters, the values of which serve to complete the definition of the function. We may have a set of quantities $w$, $t$, \ldots\ each of them a function of the same variables $u, v, \ldots$; and this relation may be expressed by means of a single functional symbol $\phi$, $(w, t, \ldots) = \phi(u, v, \ldots)$; but, as to this, more hereafter.

The notion of a function is applicable in geometry and mechanics as well as in analysis; for instance, a point $Q$, the position of which depends upon that of a variable point $P$, may be regarded as a function of the point $P$: but here, substituting for the points themselves the coordinates (of any kind whatever) which determine their positions, we may say that the coordinates of $Q$ are each of them a function of the coordinates of $P$, and we thus return to the analytical notion of a function. And in what follows a function is regarded exclusively in this point of view, viz.\ the variables are regarded as numbers; and we attend to the case of a function of one variable $v = f u$. But it has been remarked (see Equation) that it is not allowable to confine the attention to real numbers; a number $u$ must in general be taken to be a complex number $u = x + iy$, $x$ and $y$ being real numbers, each susceptible of continuous variation between the limits $-\infty$, $+\infty$, and $i$ denoting $\sqrt{-1}$. In regard to any particular function, $fu$, although it may for some purposes be sufficient to know the value of the function for any real value whatever of $u$, yet to attend only to the real values of $u$ is an essentially incomplete view of the question: to properly know the function, it is necessary to consider $u$ under the aforesaid imaginary or complex form $u = x + iy$.

To a given value $x + iy$ of $u$ there corresponds in general for $v$ a value or values of the like form $v = x' + iy'$, and we obtain a geometrical notion of the meaning of the functional relation $v = fu$ by regarding $x, y$ as rectangular coordinates of a point $P$ in a plane $\Pi$, and $x', y'$ as rectangular coordinates of a point $P'$ in a plane (for greater convenience a different plane) $\Pi'$; $P, P'$ are thus the geometrical representations, or representative points, of the variables $u = x + iy$ and $v = x' + iy'$ respectively: and, according to a locution above referred to, the point $P'$ might be regarded as a function of the point $P$: a given value of $u = x + iy$ is thus represented by a point $P$ in the plane $\Pi$, and corresponding hereto we have a point or points $P'$ in the plane $\Pi'$, representing (if more than one, each of them) a value of the variable $v = x' + iy'$. And, if we attend only to the values of $u$ as corresponding to a series of positions of the representative point $P$, we have the notion of the ``path'' of a complex variable $u = x + iy$.

\smallskip
\begin{center}\textbf{Known Functions.}\end{center}

\subsection*{\S\,1.}

The most simple kind of function is the rational and integral function. We have the series of powers $u^{2}$, $u^{3}$, \ldots\ each calculable not only for a real but also for a complex value of $u$, $(x + iy)^{2} = x^{2} - y^{2} + 2ixy$, $(x + iy)^{3} = x^{3} - 3xy^{2} + i(3x^{2}y - y^{3})$, \&c., and thence, $a, b, \ldots$ being real or complex numbers, the general form $a + bu + cu^{2} + \ldots + ku^{m}$, of a rational and integral function of the order $m$. And taking two such functions, say of the orders $m$ and $n$ respectively, the quotient of one of these by the other represents the general form of a rational function of $u$.

The function which next presents itself is the algebraical function, and in particular the algebraical function expressible by radicals. To take the most simple case, suppose ($m$ being a positive integer) that $v^{m} = u$; $v$ is here the irrational function $= u^{\frac{1}{m}}$. Obviously, if $u$ is real and positive, there is always a real and positive value of $v$, calculable to any extent of approximation from the equation $v^{m} = u$, which serves as the definition of $u^{\frac{1}{m}}$; but it is known (see Equation) that, as well in this case as in the general case where $u$ is a complex number, there are in fact $m$ values of the function $u^{\frac{1}{m}}$; and that for their determination we require the theory of the so-called circular functions sine and cosine; and these depend on the exponential function $\exp u$, or, as it is commonly written, $e^{u}$, which has for its inverse the logarithmic function $\log u$; these are all of them transcendental functions.

\subsection*{\S\,2.}

In a rational and integral function $a + bu + cu^{2} + \ldots + ku^{m}$, the number of terms is finite, and the coefficients $a, b, \ldots, k$ may have any values whatever, but if we imagine a like series $a + bu + cu^{2} + \ldots$ extending to infinity, \textit{non constat} that such an expression has any calculable value, that is, any meaning at all: the coefficients $a, b, c, \ldots$ must be such as, either for every value whatever of $u$ (that is, for every finite value) or for values included within certain limits, to make the series convergent. It is easy to see that the values of $a, b, c, \ldots$ may be such as to make the series always convergent; for instance, this is the case for the exponential function,
\[
\exp u = 1 + u + \frac{u^{2}}{1 \cdot 2} + \frac{u^{3}}{1 \cdot 2 \cdot 3} + \&\text{c.};
\]
taking for the moment $u$ to be real and positive, then it is evident that however large $u$ may be, the successive terms will become ultimately smaller and smaller, and the series will have a determinate calculable value. A function thus expressed by means of a convergent infinite series is not in general algebraical, and when it is not so, it is said to be transcendental (but observe that it is in nowise true that we have thus the most general form of a transcendental function): in particular, the exponential function above written down is not an algebraical function.

By forming the expression of $\exp v$, and multiplying together the two series, we derive the fundamental property
\[
\exp u \cdot \exp v = \exp(u + v);
\]
whence also
\[
\exp x \cdot \exp iy = \exp(x + iy),
\]
so that $\exp(x + iy)$ is given as the product of the two series $\exp x$ and $\exp iy$. As regards this last, if in place of $u$ we actually write the value $iy$, we find
\[
\exp iy = \left(1 - \frac{y^{2}}{1 \cdot 2} + \frac{y^{4}}{1 \cdot 2 \cdot 3 \cdot 4} - \ldots\right) + i\left(y - \frac{y^{3}}{1 \cdot 2 \cdot 3} + \ldots\right),
\]
where obviously each series is convergent and actually calculable for any real value whatever of $y$. Calling the two series cosine $y$ and sine $y$ respectively, or in the ordinary abbreviated notation $\cos y$ and $\sin y$, the equation is
\[
\exp iy = \cos y + i \sin y;
\]
and if we herein for $y$ write $z$, and multiply the two expressions together, observing that the product will be $= \exp i(y + z)$, we obtain the fundamental equations
\begin{align*}
\cos(y + z) &= \cos y \cos z - \sin y \sin z, \\
\sin(y + z) &= \sin y \cos z + \sin z \cos y,
\end{align*}
for the functions sine and cosine.

\subsection*{\S\,3.}

Taking $y$ as an angle, and defining as usual the sine and cosine as the ratios of the perpendicular and base respectively to the radius, the sine and cosine will be functions of $y$; and we obtain geometrically the foregoing fundamental equations for the sine and cosine: but in order to the truth of the foregoing equation $\exp iy = \cos y + i\sin y$, it is further necessary that the angle should be measured in circular measure, that is, by the ratio of the arc to the radius; so that $\pi$ denoting as usual the number $3.14159\ldots$, the measure of a right angle is $= \frac{1}{2}\pi$. And this being so, the functions sine and cosine, obtained as above by consideration of the exponential function, have their ordinary geometrical significations.

\subsection*{\S\,4.}

The foregoing investigation was given in detail in order to the completion of the theory of the irrational function $u^{\frac{1}{m}}$. We henceforth take the theory of the circular functions as known, and speak of $\tan u$, \&c., as the occasion may arise.

We have
\[
x + iy = r(\cos\theta + i\sin\theta),
\]
where, writing $\sqrt{x^{2} + y^{2}}$ to denote the positive value of the square root, we have
\[
r = \sqrt{x^{2} + y^{2}}, \quad \cos\theta = \frac{x}{\sqrt{x^{2} + y^{2}}}, \quad \sin\theta = \frac{y}{\sqrt{x^{2} + y^{2}}},
\]
and therefore also $\tan\theta = \frac{y}{x}$.

Treating $x, y$ as the rectangular coordinates of a point $P$, $r$ is the distance (regarded as positive) of this point from the origin, and $\theta$ is the inclination of $r$ to the positive part of the axis of $x$; to fix the ideas $\theta$ may be regarded as lying within the limits $0$, $\pi$, or $0$, $-\pi$, according as $y$ is positive or negative; $\theta$ is thus completely determinate, except in the case, $x$ negative, $y = 0$, for which $\theta$ is $= \pi$ or $-\pi$ indifferently.

And if $u = x + iy$, we hence have
\[
(x + iy)^{\frac{1}{m}} = r^{\frac{1}{m}}\left(\cos\frac{\theta + 2s\pi}{m} + i\sin\frac{\theta + 2s\pi}{m}\right),
\]
where $r^{\frac{1}{m}}$ is real and positive and $s$ has any positive or negative integer value whatever: but we thus obtain for $u^{\frac{1}{m}}$ only the $m$ values corresponding to the values $0, 1, 2, \ldots, m - 1$ of $s$. More generally we may, instead of the index $\frac{1}{m}$, take the index to be any rational fraction $\frac{n}{m}$. Supposing this to be in its least terms, and $m$ to be positive, the number of distinct values is always $= m$. If instead of $\frac{n}{m}$ we take the index to be the general real or complex quantity $\mu$, we have $u^{\mu}$, no longer an algebraical function of $u$, and having in general an infinity of values.

\subsection*{\S\,5.}

The foregoing equation $\exp(x + y) = \exp x \cdot \exp y$ is, in fact, the equation of indices, $a^{x+y} = a^{x} \cdot a^{y}$; $\exp x$ is thus the same thing as $e^{x}$, where $e$ denotes a properly determined number, and putting $e^{x}$ equal to the series, and then writing $x = 1$, we have $e = 1 + 1 + \frac{1}{1 \cdot 2} + \frac{1}{1 \cdot 2 \cdot 3} + \&$c., that is, $e = 2.71828\ldots$ But as well theoretically as for convenience of printing, there is considerable advantage in the use of the notation $\exp u$.

From the equation, $\exp iy = \cos y + i\sin y$, we deduce $\exp(-iy) = \cos y - i\sin y$, and thence
\begin{align*}
\cos y &= \tfrac{1}{2}\{\exp(iy) + \exp(-iy)\}, \\
\sin y &= \tfrac{1}{2i}\{\exp(iy) - \exp(-iy)\};
\end{align*}
if we write herein $ix$ instead of $y$ we have
\begin{align*}
\cos ix &= \tfrac{1}{2}\{\exp x + \exp(-x)\}, \\
\sin ix &= \tfrac{1}{2i}\{\exp x - \exp(-x)\},
\end{align*}
viz.\ these values are
\begin{align*}
\cos ix &= 1 + \frac{x^{2}}{1 \cdot 2} + \frac{x^{4}}{1 \cdot 2 \cdot 3 \cdot 4} + \ldots, \\
\sin ix &= \frac{x}{1} + \frac{x^{3}}{1 \cdot 2 \cdot 3} + \ldots,
\end{align*}
each of them real when $x$ is real. The functions in question $1 + \frac{x^{2}}{1 \cdot 2} + \ldots$ and $x + \frac{x^{3}}{1 \cdot 2 \cdot 3} + \ldots$, regarded as functions of $x$, are termed the hyperbolic cosine and sine, and are represented by the notations $\cosh x$ and $\sinh x$ respectively; and similarly we have the hyperbolic tangent $\tanh x$, \&c.: although it is easy to remember that $\cos ix$, $\frac{1}{i}\sin ix$ are, in fact, real functions of $x$, and to understand accordingly the formulae wherein they occur, yet the use of these notations of the hyperbolic functions is often convenient.

\subsection*{\S\,6.}

Writing $u = \exp v$, then $v$ is conversely a function of $u$ which is called the logarithm (hyperbolic logarithm, to distinguish it from the tabular or Briggian logarithm), and we write $v = \log u$, or what is the same thing, we have $u = \exp(\log u)$: and it is clear that if $u$ be real and positive there is always a real and positive value of $\log u$, in particular the real logarithm of $e$ is $= 1$; it is however to be observed that the logarithm is not a one-valued function, but has an infinity of values corresponding to the different integer values of a constant $s$; in fact, if $\log u$ be any one of its values, then $\log u + 2s\pi i$ is also a value, for we have $\exp(\log u + 2s\pi i) = \exp\log u \cdot \exp 2s\pi i$, or since $\exp 2s\pi i$ is $= 1$, this is $= u$; that is, $\log u + 2s\pi i$ is a value of the logarithm of $u$.

We have
\[
uv = \exp(\log uv) = \exp\log u \cdot \exp\log v,
\]
and hence the equation which is commonly written
\[
\log uv = \log u + \log v,
\]
but which requires the addition on one side of a term $2s\pi i$. And reverting to the equation $x + iy = r(\cos\theta + i\sin\theta)$, or as it is convenient to write it, $x + iy = r\exp i\theta$, we hence have
\[
\log(x + iy) = \log r + i(\theta + 2s\pi),
\]
where $\log r$ may be taken to denote the real logarithm of the real positive quantity $r$, and $\theta$ the completely determinate angle defined as already mentioned.

Reverting to the function $u^{\frac{1}{m}}$, we have $u = \exp\log u$, and thence $u^{\frac{1}{m}} = \exp(\frac{1}{m}\log u)$, which, on account of the infinity of values of $\log u$, has in general (as before remarked) an infinity of values; if $u = e$, then $e^{\frac{1}{m}} = \exp(\frac{1}{m}\log e)$, has in general in like manner an infinity of values, but in regarding $e^{\frac{1}{m}}$ as identical with the one-valued function $\exp\frac{1}{m}$, we take $\log e$ to be $=$ its real value, 1.

The inverse functions $\cos^{-1}x$, $\sin^{-1}x$, $\tan^{-1}x$, are in fact logarithmic functions; thus in the equation $\exp ix = \cos x + i\sin x$, writing first $\cos x = u$, the equation becomes $\exp(i\cos^{-1}u) = u + \sqrt{u^{2} - 1}$, or we have $\cos^{-1}u = \frac{1}{i}\log(u + \sqrt{u^{2} - 1})$, and from the same equation, writing secondly $\sin x = u$, we have $\sin^{-1}u = \frac{1}{i}\log(\sqrt{1 - u^{2}} + iu)$. But the formula for $\tan^{-1}u$ is a more elegant one, as not involving the radical $\sqrt{1 - u^{2}}$; we have
\[
\tan x = \frac{1}{i}\frac{\exp ix - \exp(-ix)}{\exp ix + \exp(-ix)} = \frac{1}{i}\frac{\exp 2ix - 1}{\exp 2ix + 1},
\]
and thence $\exp 2ix = \dfrac{1 + i\tan x}{1 - i\tan x}$; that is,
\[
x = \frac{1}{2i}\log\frac{1 + i\tan x}{1 - i\tan x},
\]
or, if $\tan x = u$, then
\[
\tan^{-1}u = \frac{1}{2i}\log\frac{1 + iu}{1 - iu}.
\]

The logarithm (or inverse exponential function) and the inverse circular functions present themselves as the integrals of algebraic functions:
\[
\int\frac{dx}{x} = \log x,
\]
whence also $\displaystyle\int\frac{dx}{1 + x^{2}} = \frac{1}{2i}\log\frac{1 + ix}{1 - ix} = \tan^{-1}x$, and $\displaystyle\int\frac{dx}{\sqrt{1 - x^{2}}} = \sin^{-1}x$.

\subsection*{\S\,7.}

Each of the functions $\exp u$, $\sin u$, $\cos u$, $\tan u$, \&c., is a one-valued function of $u$: in this respect analogous to a rational function of $u$; and there are further analogies of $\exp u$, $\sin u$, $\cos u$, to a rational and integral function; and of $\tan u$, $\sec u$, \&c., to a rational non-integral function.

A rational and integral function has a certain number of roots, or zeros each of a given multiplicity, and is completely determined (except as to a constant factor) when the several roots and the multiplicity of each of them is given; i.e., if $a, b, c, \ldots$ be the roots, $p, q, r, \ldots$ their multiplicities, then the form is $a(x - a)^{p}(x - b)^{q}\ldots$; a rational (non-integral) function has a certain number of infinities, or ``poles,'' each of them of a given multiplicity; viz.\ the infinities are the roots or zeros of the rational and integral function which is its denominator.

The function $\exp u$ has no finite roots, but an infinity of roots each $= \infty$; this appears from the equation $\exp u = \left(1 + \frac{u}{n}\right)^{n}$, where $n$ is indefinitely large and positive. The function $\sin u$ has the roots $s\pi$, where $s$ is any positive or negative integer, zero included; or, what is the same thing, its roots are $0$ and $\pm s\pi$, $s$ now denoting any positive integer from 1 to $\infty$; each of these is a simple root, and we in fact have
\[
\sin u = u\prod_{s}\left(1 - \frac{u^{2}}{s^{2}\pi^{2}}\right).
\]
Similarly the roots of $\cos u$ are $(s + \frac{1}{2})\pi$, $s$ denoting any positive or negative integer, zero included, or, what is the same thing, they are $\pm(s + \frac{1}{2})\pi$, $s$ now denoting any positive integer from 1 to $\infty$; each root is simple, and we have
\[
\cos u = \prod_{s}\left(1 - \frac{u^{2}}{(s + \frac{1}{2})^{2}\pi^{2}}\right).
\]
Obviously $\tan u$, as the quotient $\sin u \div \cos u$, has both roots and infinities, its roots being the roots of $\sin u$, its infinities the roots of $\cos u$; $\sec u$, as the reciprocal of $\cos u$, has infinities only, these being the roots of $\cos u$, \&c.

\subsection*{\S\,8.}

In the foregoing expression $\sin u = u\prod_{s}(1 - \frac{u^{2}}{s^{2}\pi^{2}})$, the product must be understood to mean the limit of $\prod_{-n}^{n}(1 - \frac{u^{2}}{s^{2}\pi^{2}})$ for an indefinitely large positive integer value of $n$, viz.\ the product is first to be formed for the values $s = 1, 2, 3, \ldots$ up to a determinate number $n$, and then $n$ is to be taken indefinitely large. If, separating the positive and the negative values of $s$, we consider the product $u\prod_{1}^{m}(1 + \frac{u}{s\pi}) \cdot \prod_{1}^{n}(1 - \frac{u}{s\pi})$, where in the first product $s$ has all the positive integer values from 1 to $m$, and in the second product $s$ has all the positive integer values from 1 to $n$, then by making $m$ and $n$ each of them indefinitely large, the function does not approximate to $\sin u$ unless $m : n$ be a ratio of equality\textsuperscript{*}.

\smallskip
\small\noindent\rule{0.4\textwidth}{0.4pt}\\
\textsuperscript{*} The value of the function in question $u\prod_{1}^{m}(1 + \frac{u}{s\pi}) \cdot \prod_{1}^{n}(1 - \frac{u}{s\pi})$, when $m, n$ are each indefinitely large, but $\frac{m}{n}$ not $= 1$, is $= \left(\frac{m}{n}\right)^{\frac{u}{\pi}}\sin u$.
\normalsize

\subsection*{\S\,9.}

The functions $\sin u$, $\cos u$, are periodic, having the period $2\pi$, $\cos(u + 2\pi) = \cos u$; and the half-period $\pi$, $\cos(u + \pi) = -\cos u$; the periodicity may be verified by means of the foregoing fractional forms, but some attention is required; thus writing, as we may do, $\sin u = \frac{u\prod(u \pm s\pi)}{\ldots}$, where $s$ extends from $-n$ to $n$, $n$ ultimately infinite, if for $u$ we write $u + \pi$, each factor of the numerator is changed into the following one and the numerator is unaltered, save only that there is an introduced factor $u + (n + 1)\pi$ at the superior limit, and an omitted factor $u - n\pi$ at the inferior limit; the ratio of these, $(u + n + 1\pi) \div (u - n\pi)$, for $n$ infinite is $= -1$, and we thus have, as we should have, $\sin(u + \pi) = -\sin u$.

The most general periodic function having no infinities, and each root a simple root, and having a given period $a$, has the form $A\sin\frac{2\pi u}{a} + B\cos\frac{2\pi u}{a}$, or, what is the same thing, $C\sin\left(\frac{2\pi u}{a} + D\right)$.

\subsection*{\S\,10.}

We come now to the Elliptic Functions. These arose from the consideration of the integral $\int\frac{dx}{\sqrt{X}}$, where $R$ is a rational function of $x$, and $X$ is the general rational and integral quartic function $\alpha x^{4} + \beta x^{3} + \gamma x^{2} + \delta x + \epsilon$; a form arrived at was
\[
\int\frac{dx}{\sqrt{1 - x^{2} \cdot 1 - k^{2}x^{2}}} = \int\frac{d\phi}{\sqrt{1 - k^{2}\sin^{2}\phi}},
\]
on putting therein $x = \sin\phi$, and this last integral was represented by $F\phi$, and called the elliptic integral of the first kind. In the particular case $k = 0$, the integral is $= \sin^{-1}x$, and it thus appears that $F\phi$ is of the nature of an inverse function: for passing to the direct functions we write $F\phi = u$, and consider $\phi$ as hereby determined as a function of $u$, $\phi = \text{amplitude of } u$, or for shortness $\am u$. And the functions $\sin\phi$, $\cos\phi$, and $\sqrt{1 - k^{2}\sin^{2}\phi}$ were then considered as functions of the amplitude, and written $\sin\am u$, $\cos\am u$, $\Delta\am u$: these were afterwards written $\sn u$, $\cn u$, $\dn u$, which may be regarded either as mere abbreviations of the former functional symbols, or (in a different point of view) as functions, no longer of $\am u$, but of $u$ itself as the argument of the functions: $\sn$ is thus a function in some respects analogous to a sine, and $\cn$ and $\dn$ functions analogous to a cosine; they have the corresponding property that the three functions of $u + v$ are expressible in terms of the functions of $u$ and of $v$. The following formulae may be mentioned:
\begin{align*}
\cn^{2}u &= 1 - \sn^{2}u, & \dn^{2}u &= 1 - k^{2}\sn^{2}u, \\
\sn' u &= \cn u \dn u, & \cn' u &= -\sn u \dn u, & \dn' u &= -k^{2}\sn u \cn u,
\end{align*}
where the accent denotes differentiation in regard to $u$; and the addition-formulae
\begin{align*}
\sn(u + v) &= \sn u \cn v \dn v + \sn v \cn u \dn u, \\
\cn(u + v) &= \cn u \cn v - \sn u \dn u \sn v \dn v, \\
\dn(u + v) &= \dn u \dn v - k^{2}\sn u \cn u \sn v \cn v,
\end{align*}
each of the expressions on the right-hand side being the numerator of a fraction of which
\[
\text{Denom.} = 1 - k^{2}\sn^{2}u \sn^{2}v.
\]
It may be remarked that any one of the fractional expressions, differentiated in regard to $u$ and to $v$ respectively, gives the same result; such expression is therefore a function of $u + v$, and the addition-formulae can be thus directly verified.

\subsection*{\S\,11.}

The existence of a denominator in the addition-formulae suggests that $\sn$, $\cn$, $\dn$ are not, like the sine and cosine, functions having zeros only, without infinities; they are in fact functions, having each its own zeros, but having a common set of infinities; moreover, the zeros and the infinities are simple zeros and infinities respectively. And this further suggests, what in fact is the case, that the three functions are quotients having each its own numerator but a common denominator, say they are the quotients of four $\vartheta$-functions, each of them having zeros only (and these simple zeros) but no infinities.

The functions $\sn$, $\cn$, $\dn$, but not the $\vartheta$-functions, are moreover doubly periodic; that is, there exist values $2\omega$, $2\nu$, $= 4K$ and $4i(K + iK')$ in the ordinary notation, such that the $\sn$, $\cn$, or $\dn$ of $u + 2\omega$, and the $\sn$, $\cn$, and $\dn$ of $u + 2\nu$ are equal to the $\sn$, $\cn$, and $\dn$ respectively of $u$; or say that $\phi(u + 2\omega) = \phi(u + 2\nu) = \phi u$, where $\phi$ is any one of the three functions.

As regards this double periodicity, it is to be observed that the equations $\phi(u + 2\omega) = \phi u$, $\phi(u + 2\nu) = \phi u$, imply $\phi(u + 2m\omega + 2n\nu) = \phi u$, and hence it easily follows that if $\omega$, $\nu$ were commensurable, say if they were multiples of some quantity $\sigma$, we should have $\phi(u + 2\sigma) = \phi u$, an equation which would replace the original two equations $\phi(u + 2\omega) = \phi u$, $\phi(u + 2\nu) = \phi u$, or there would in this case be only the single period $\sigma$; $\omega$ and $\nu$ must therefore be incommensurable. And this being so, they cannot have a real ratio, for if they had, the integer values $m, n$ could be taken such as to make $2m\omega + 2n\nu = k$ times a given real or imaginary value, $k$ as small as we please; the ratio $\omega : \nu$ must be therefore imaginary, as is in fact the case when the values are $4K$ and $4i(K + iK')$.

\subsection*{\S\,12.}

The function $\sn u$ has the zero $u$ and the zeros $m\omega + n\nu$, $m$ and $n$ any positive or negative integers whatever; and this suggests that the numerator of $\sn u$ is equal to a doubly infinite product, (Cayley, ``On the Inverse Elliptic Functions,'' \textit{Camb.\ Math.\ Jour.}\ t.~iv., 1845, [24]; and ``M\'{e}moire sur les fonctions doublement p\'{e}riodiques,'' \textit{Liouville}, t.~x., 1845, [25]). The numerator is equal to
\[
u \prod_{m,n}\left(1 - \frac{u}{m\omega + n\nu}\right),
\]
$m$ and $n$ having any positive or negative integer values whatever, including zero, except that $m, n$ must not be simultaneously $= 0$, these values being taken account of in the factor $u$ outside the product. But until further defined, such a product has no definite value, and consequently no meaning whatever. Imagine $m, n$ to be coordinates, and suppose that we have, surrounding the origin, a closed curve having the origin for its centre, i.e.\ the curve is such that, if $\alpha, \beta$ be the coordinates of a point thereof, then $-\alpha, -\beta$ are also the coordinates of a point thereof; suppose further that the form of the curve is given, but that its magnitude depends upon a parameter $h$, and that the curve is such that, when $h$ is indefinitely large, each point of the curve is at an indefinitely large distance from the origin; for instance, the curve might be a circle or ellipse, or a parallelogram, the origin being in each case the centre. Then if in the double product, taking the value of $h$ as given, we first give to $m, n$ all the positive or negative integer values (the simultaneous values $0, 0$ excluded) which correspond to points within the curve, and then make $h$ indefinitely large, the product will thus have a definite value: but this value will still be dependent on the form of the curve.

Moreover, varying in any manner the form of the curve, the ratio of the two values of the double product will be $= \exp\beta u^{2}$, where $\beta$ is a determinate value depending only on the forms of the two curves; or, what is the same thing, if we first give to the curve a certain form, say we take it to be a circle, and then give it any other form, the product in the latter case is equal to its former value multiplied by $\exp\beta u^{2}$, where $\beta$ depends only upon the form of the curve in the latter case.

Considering the form of the bounding curve as given, and writing the double product in the form
\[
u \prod_{m,n}\left(1 - \frac{u}{m\omega + n\nu}\right),
\]
the simultaneous values $m = 0, n = 0$ being now admitted in the numerator, although still excluded from the denominator, then if we write for instance $u + 2\omega$ instead of $u$, each factor in the numerator is changed into a contiguous factor, and the numerator remains unaltered, except that we introduce certain factors which lie outside the bounding curve, and omit certain factors which lie inside the bounding curve; we, in fact, affect the result by a singly infinite series of factors belonging to points adjacent to the bounding curve: and it appears on investigation that we thus introduce a constant factor $\exp\gamma(u + \omega)$. The final result thus is that the product
\[
u \prod\left(1 - \frac{u}{m\omega + n\nu}\right)
\]
does not remain unaltered when $u$ is changed into $u + 2\omega$, but that it becomes therefore affected with a constant factor, $\exp\gamma(u + \omega)$. And similarly the function does not remain unaltered when $u$ is changed into $u + 2\nu$, but it becomes affected with a factor, $\exp\delta(u + \nu)$. The bounding curve may however be taken such that the function is unaltered when $u$ is changed into $u + 2\omega$: this will be the case if the curve is a rectangle such that the length in the direction of the axis of $m$ is infinitely great in comparison of that in the direction of the axis of $n$; or it may be taken such that the function is unaltered when $u$ is changed into $u + 2\nu$: this will be so if the curve be a rectangle such that the length in the direction of the axis of $n$ is indefinitely great in comparison with that in the direction of the axis of $m$; but the two conditions cannot be satisfied simultaneously.

\subsection*{\S\,13.}

We have three other like functions, viz.\ writing for shortness $\bar{m}$, $\bar{n}$ to denote $m + \frac{1}{2}$, $n + \frac{1}{2}$ respectively, and $(m, n)$ to denote $m\omega + n\nu$, then the four functions are
\begin{multline*}
\vartheta_{1} = \nu \prod_{m,n}\left(1 - \frac{u}{(m, n)}\right), \quad
\vartheta_{2} = \nu \prod_{m,n}\left(1 - \frac{u + \frac{1}{2}\omega}{(m, n)}\right), \\
\vartheta_{3} = \nu \prod_{m,n}\left(1 - \frac{u + \frac{1}{2}\nu}{(m, n)}\right), \quad
\vartheta_{4} = \nu \prod_{m,n}\left(1 - \frac{u + \frac{1}{2}\omega + \frac{1}{2}\nu}{(m, n)}\right),
\end{multline*}
the bounding curve being in each case the same; and, dividing the first three of these each by the last, we have (except as to constant factors) the three functions $\sn u$, $\cn u$, $\dn u$; writing in each of the four functions $u + 2\omega$ or $u + 2\nu$ in place of $u$, the functions acquire each of them the same exponential factor $\exp\gamma(u + \omega)$, or $\exp\delta(u + \nu)$, and the quotient of any two of them, and therefore each of the functions $\sn u$, $\cn u$, $\dn u$, remains unaltered.

It is easily seen that, disregarding constant factors, the four $\vartheta$-functions are in fact one and the same function with different arguments, or they may be written $\vartheta u$, $\vartheta(u + \frac{1}{2}\omega)$, $\vartheta(u + \frac{1}{2}\nu)$, $\vartheta(u + \frac{1}{2}\omega + \frac{1}{2}\nu)$; by what precedes, the functions may be so determined that they shall remain unaltered when $u$ is changed into $u + 2\omega$, that is, be singly periodic, but that the change $u$ into $u + 2\nu$ shall affect them each with the same exponential factor $\exp\delta(u + \nu)$.

\subsection*{\S\,14.}

Taking the last-mentioned property as a definition of the function $\vartheta$, it appears that $\vartheta u$ may be expressed as a sum of exponentials
\[
\vartheta u = \sum_{m} A\exp\frac{-}{}(\nu m^{2} + um),
\]
where the summation extends to all positive and negative integer values of $m$, including zero. In fact, if we first write herein $u + 2\omega$ instead of $u$, then in each term the index of the exponential is altered by $-2\omega m$, $= 2\nu m i$, and the term itself thus remains unaltered; that is, $\vartheta(u + 2\omega) = \vartheta u$. But writing $u + 2\nu$ in place of $u$, each term is changed into the succeeding term, multiplied by the factor $\exp\frac{-}{}(m + 1)\nu$; in fact, making the change in question $u$ into $u + 2\nu$, and writing also $m - 1$ in place of $m$, $\nu m^{2} + um$ becomes $\nu(m - 1)^{2} + (u + 2\nu)(m - 1)$, $= \nu m^{2} + um - u - \nu$, and we thus have
\[
\vartheta(u + 2\nu) = \exp\frac{-}{}(u + \nu) \cdot \vartheta u.
\]
In order to the convergency of the series it is necessary that $\exp\frac{-}{}\nu m^{2}$ should vanish for indefinitely large values of $m$, and this will be so if $\frac{-}{}\nu$ be a complex quantity of the form $\alpha + \beta i$, $\alpha$ negative; for instance, this will be the case if $\omega$ be real and positive and $\nu$ be $= i$ multiplied by a real and positive quantity. The original definition of $\vartheta$ as a double product seems to put more clearly in evidence the real nature of this function, but the new definition has the advantage that it admits of extension to the $\vartheta$-functions of two or more variables.

The elliptic functions $\sn u$, $\cn u$, $\dn u$, have thus been expressed each of them as the quotient of two $\vartheta$-functions, but the question arises to express conversely a $\vartheta$-function by means of the elliptic functions; the form is found to be
\[
\vartheta u = C\exp\left(Au^{2} + B\int\sn^{2}u\,du\right),
\]
viz.\ $\vartheta u$ is expressible as an exponential, the index of which depends on the double integral $\int\sn^{2}u\,du$.

The object has been to explain the general nature of the elliptic functions $\sn u$, $\cn u$, $\dn u$, and of the $\vartheta$-functions with which they are thus intimately connected; it would be out of place to go into the theories of the multiplication, division, and transformation of the elliptic functions, or into the theory of the elliptic integrals, and of the application of the $\vartheta$-functions to the representation of the elliptic integrals of the second and third kinds.

\subsection*{\S\,15.}

The reasoning which shows that for a doubly periodic function the ratio of the two periods $2\omega$, $2\nu$ is imaginary shows that we cannot have a function of a single variable, which shall be triply periodic, or of any higher order of periodicity. For if the periods of a triply periodic function $\phi(u)$ were $2\omega$, $2\nu$, $2\chi$, then $m, n, p$ being any positive or negative integer values, we should have $\phi(u + 2m\omega + 2n\nu + 2p\chi) = \phi u$; $\omega, \nu, \chi$ must be incommensurable, for if not, the three periods would really reduce themselves to two periods, or to a single period; and being incommensurable, it would be possible to determine the integers $m, n, p$ in such manner that the real part and also the coefficient of $i$ of the expression $m\omega + n\nu + p\chi$ shall be each of them as small as we please, say $< \epsilon$; then $\phi(u + \epsilon) = \phi u$, and thence $\phi(u + k\epsilon) = \phi u$ ($k$ an integer), and $k\epsilon$ as near as we please to any given real or imaginary value whatever. We have thus the nugatory result $\phi u = \text{a constant}$, or at least the function if not a constant is a function of an infinitely and perpetually discontinuous kind, a conception of which can hardly be formed. But a function of two variables may be triply or quadruply periodic: viz.\ we may have a function $\phi(u, v)$ having for $u, v$ the simultaneous periods $2\omega$, $2\omega'$; $2\nu$, $2\nu'$; $2\chi$, $2\chi'$; or, what is the same thing, it may be such that, $m, n, p, q$ being any integers whatever, we have
\[
\phi(u + 2m\omega + 2n\nu + 2p\chi + 2q\psi,\; v + 2m\omega' + 2n\nu' + 2p\chi' + 2q\psi') = \phi(u, v);
\]
and similarly a function of $2n$ variables may be $2n$-tuply periodic.

It is, in fact, in this manner that we pass from the elliptic functions and the single $\vartheta$-functions to the hyperelliptic or Abelian functions and the multiple $\vartheta$-functions; the case next succeeding the elliptic functions is when we have $X$, $Y$ the same rational and integral sextic functions of $x$, $y$ respectively, and then writing
\[
\int\frac{dx}{\sqrt{X}} + \int\frac{dy}{\sqrt{Y}} = du, \quad
\int\frac{x\,dx}{\sqrt{X}} + \int\frac{y\,dy}{\sqrt{Y}} = dv,
\]
we regard certain symmetrical functions of $x, y$, in fact, the ratios of ($2^{4} =$) 16 such symmetrical functions, as functions of $(u, v)$; say we thus have 15 hyperelliptic functions $f(u, v)$, analogous to the 3 elliptic functions $\sn u$, $\cn u$, $\dn u$, and being quadruply periodic. And these are the quotients of 16 double $\vartheta$-functions $\vartheta(u, v)$, the general form being
\[
\vartheta(u, v) = \sum\sum A\exp\frac{-}{}(a, b, b)(m, n)^{2} + mu + nv,
\]
where the summations extend to all positive and negative integer values of $(m, n)$: and we thus see the form of the $\vartheta$-function for any number of variables whatever. The epithet ``hyperelliptic'' is used in the case where the differentials are of the form just mentioned $\frac{dx}{\sqrt{X}}$, where $X$ is a rational and integral function of $x$; the epithet ``Abelian'' extends to the more general case where the differential involves the irrational function of $x$, determined by any rational and integral equation $\phi(x, y) = 0$ whatever.

As regards the literature of the subject, it may be noticed that the various memoirs by Riemann, 1851--1866, are republished in the collected edition of his works, Leipsic, 1876; and shortly after his death we have the \textit{Theorie der Abel'schen Functionen}, by Clebsch and Gordan, Leipsic, 1866. Preceding this, we have by MM.\ Briot and Bouquet, the \textit{Th\'{e}orie des Fonctions doublement p\'{e}riodiques et en particulier des Fonctions Elliptiques}, Paris, 1859, the results of which are reproduced and developed in their larger work, \textit{Th\'{e}orie des Fonctions Elliptiques}, 2nd ed., Paris, 1875.

\subsection*{\S\,16.}

It is proper to mention the gamma ($\Gamma$) or $\Pi$ function, $\Gamma(n + 1) = \Pi n = 1 \cdot 2 \cdot 3 \ldots n$, when $n$ is a positive integer. In the case just referred to, $n$ a positive integer, this presents itself almost everywhere in analysis,---for instance, the binomial coefficients, and the coefficients of the exponential series are expressible by means of such functions of a number $n$. The definition for any real positive value of $n$ is taken to be
\[
\Pi n = \int_{0}^{\infty} x^{n}\exp(-x)\,dx;
\]
it is then shown that, $n$ being real and positive, $\Gamma(n + 1) = n\Gamma n$, and by assuming that this equation holds good for positive or negative real values of $n$, the definition is extended to real negative values; the equation gives $\Pi 1 = 0 \cdot \Pi 0$, that is, $\Pi 0 = \infty$; and similarly $\Pi(-n) = \infty$, where $-n$ is any negative integer. The definition by the definite integral has been or may be extended to imaginary values of $n$, but the theory is not an established one. A definition extending to all values of $n$ is that of Gauss
\[
\Pi n = \text{limit } \frac{k^{n} \cdot k!}{(n + 1)(n + 2)(n + 3)\ldots(n + k)},
\]
the ultimate value of $k$ being $= \infty$; but the function is chiefly considered for real values of the variable.

A formula for the calculation, when $x$ has a large real and positive value, is
\[
\Pi x = \sqrt{2\pi}\; x^{x + \frac{1}{2}}\exp(-x + \frac{1}{12x} - \ldots),
\]
or as this may also be written, neglecting the negative powers of $x$,
\[
\Pi x = \sqrt{2\pi}\exp\{(x + \tfrac{1}{2})\log x - x\}.
\]

Another formula is $\Pi x \cdot \Pi(-x) = \dfrac{\pi x}{\sin\pi x}$: or, as this may also be written,
\[
\Pi(x - 1)\Pi(-x) = \frac{\pi}{\sin\pi x}.
\]

It is to be observed that the function $\Pi$ serves to express the product of a set of factors in arithmetical progression; we have
\[
(x + a)(x + 2a)\ldots(x + ma) = a^{m}\Pi\left(\frac{x}{a} + 1\right)\Pi\left(\frac{x}{a} + 2\right)\ldots\left(\frac{x}{a} + m\right) = a^{m}\frac{\Pi(\frac{x}{a} + m)}{\Pi\frac{x}{a}}.
\]
We can consequently express by means of it the product of any number of the factors which present themselves in the factorial expression of $\sin u$. Starting from the form
\[
\sin u = u\prod_{s}\left(1 - \frac{u^{2}}{s^{2}\pi^{2}}\right),
\]
where $\Pi$ is here as before the sign of a product of factors corresponding to the different integer values of $s$, this is thus converted into
\[
\sin u = u\Pi\left(-\frac{u}{s\pi}\right)\Pi\left(\frac{u}{s\pi}\right),
\]
or as this may also be written,
\[
\sin u = u\Pi\left(-\frac{u}{\pi}\right)\Pi\left(\frac{u}{\pi}\right)\Pi\left(-\frac{u}{2\pi}\right)\Pi\left(\frac{u}{2\pi}\right)\ldots,
\]
which, in virtue of
\[
\Pi(-\xi)\Pi\xi = \frac{\pi\xi}{\sin\pi\xi},
\]
becomes
\[
\sin u = u\Pi\left(-\frac{u}{\pi}\right)\Pi\left(\frac{u}{\pi}\right)\Pi\left(-\frac{u}{2\pi}\right)\Pi\left(\frac{u}{2\pi}\right)\ldots = u\frac{\pi u}{\sin\pi u}\ldots.
\]
Here $m$ and $n$ are large and positive; calculating the second factor by means of the formula for $\Pi x$, in this case we have the before-mentioned formula
\[
\sin u = u\prod_{s=1}^{\infty}\left(1 - \frac{u^{2}}{s^{2}\pi^{2}}\right).
\]

The gamma or $\Pi$ function is the so-called second Eulerian integral; the first Eulerian integral
\[
\int_{0}^{1}x^{p-1}(1 - x)^{q-1}\,dx = \frac{\Gamma p \cdot \Gamma q}{\Gamma(p + q)},
\]
is at once expressible in terms of $\Gamma$, and is therefore not a new function to be considered.

\subsection*{\S\,17.}

We have the function defined by its expression as a hypergeometric series
\[
F(\alpha, \beta, \gamma, u) = 1 + \frac{\alpha \cdot \beta}{1 \cdot \gamma}u + \frac{\alpha(\alpha + 1) \cdot \beta(\beta + 1)}{1 \cdot 2 \cdot \gamma(\gamma + 1)}u^{2} + \ldots,
\]
i.e.\ this expression of the function serves as a definition, if the series be finite or if, being infinite, it is convergent. The function may also be defined as a definite integral; in other words, if, in the integral
\[
\int_{0}^{1}x^{\alpha' - 1}(1 - x)^{\beta' - 1}(1 - ux)^{-\gamma'}\,dx,
\]
we expand the factor $(1 - ux)^{-\gamma'}$ in powers of $ux$, and then integrate each term separately by the formula for the second Eulerian integral, the result is
\[
\frac{\Gamma\alpha' \cdot \Gamma\beta'}{\Gamma(\alpha' + \beta')}\left\{1 + \frac{\alpha' \cdot \gamma'}{\alpha' + \beta' \cdot 1}u + \ldots\right\},
\]
which is
\[
\frac{\Gamma\alpha' \cdot \Gamma\beta'}{\Gamma(\alpha' + \beta')}\left\{1 + \frac{\alpha'}{\alpha' + \beta'}\frac{\gamma'}{1}u + \frac{\alpha'(\alpha' + 1)}{\alpha' + \beta' \cdot \alpha' + \beta' + 1}\frac{\gamma'(\gamma' + 1)}{1 \cdot 2}u^{2} + \ldots\right\},
\]
or writing $\alpha'$, $\beta'$, $\gamma'$ $= \alpha$, $\gamma - \alpha$, $\beta$ respectively, this is
\[
\frac{\Gamma\alpha \cdot \Gamma(\gamma - \alpha)}{\Gamma\gamma}F(\alpha, \beta, \gamma, u),
\]
so that the new definition is
\[
F(\alpha, \beta, \gamma, u) = \frac{\Gamma\gamma}{\Gamma\alpha \cdot \Gamma(\gamma - \alpha)}\int_{0}^{1}x^{\alpha - 1}(1 - x)^{\gamma - \alpha - 1}(1 - ux)^{-\beta}\,dx,
\]
but this is in like manner only a definition under the proper limitations as to the values of $\alpha, \beta, \gamma, u$. It is not here considered how the definition is to be extended so as to give a meaning to the function $F(\alpha, \beta, \gamma, u)$ for all values, say of the parameters $\alpha, \beta, \gamma$, and of the variable $u$. There are included a large number of special forms which are either algebraic or circular or exponential, for instance $F(\alpha, \beta, \beta, u) = (1 - u)^{-\alpha}$, \&c.; or which are special transcendents which have been separately studied, for instance, Bessel's functions, the Legendrian functions $X_{n}$ presently referred to, series occurring in the development of the reciprocal of the distance between two planets, \&c.

\subsection*{\S\,18.}

There is a class of functions depending upon a variable or variables $x, y, \ldots$ and a parameter $n$, say the function for the parameter $n$ is $X_{n}$, such that the product of two functions having the same variables, multiplied it may be by a given function of the variables, and integrated between given limits, gives a result $= 0$ or not $= 0$, according as the parameters are unequal or equal; $\int \mathfrak{A} \cdot X_{m}X_{n}\,dx = 0$, but $\int \mathfrak{A} \cdot X_{n}^{2}\,dx$ not $= 0$; the admissible values of the parameters being either any integer values, or it may be the roots of a determinate algebraical or transcendental equation; and the functions $\mathfrak{A}_{n}$ may be either algebraical or transcendental. For instance, such a function is $\cos nx$; $m$ and $n$ being integers, we have $\int_{0}^{\pi}\cos mx \cos nx\,dx = 0$, but $\int_{0}^{\pi}\cos^{2}nx\,dx = \frac{1}{2}\pi$.

Assuming the existence of the expansion of a function $fx$, in a series of multiple cosines, we thus obtain at once the well-known Fourier series, wherein the coefficient of $\cos mx$ is $= \frac{2}{\pi}\int_{0}^{\pi}\cos mx \cdot fx\,dx$. The question whether the process is applicable is elaborately discussed in Riemann's memoir (1854), \textit{Ueber die Darstellbarkeit einer Function durch eine trigonometrische Reihe}, No.~xii.\ in the collected works. And again we have the Legendrian functions, which present themselves as the coefficients of the successive powers of $a$ in the development of $(1 - 2ax + a^{2})^{-\frac{1}{2}}$, $X_{0} = 1$, $X_{1} = x$, $X_{2} = \frac{1}{2}(3x^{2} - 1)$, \&c.: here $m, n$ being any positive integers, $\int_{-1}^{1}X_{m}X_{n}\,dx = 0$, but $\int_{-1}^{1}X_{n}^{2}\,dx = \frac{2}{2n + 1}$: and we have also Laplace's functions, \&c.

\smallskip
\begin{center}\textbf{Functions in General.}\end{center}

\subsection*{\S\,19.}

In what precedes, a review has been given, not by any means an exhaustive one, but embracing the most important kinds of known functions; but there are questions to be considered in regard to functions in general.

A function of $x + iy$ has been built up by means of analytical operations performed upon $x + iy$, $(x + iy)^{2} = x^{2} - y^{2} + i \cdot 2xy$, \&c., and the question next referred to has not arisen. But observe that, knowing $x + iy$, we know $x$ and $y$, and therefore any two given functions $\phi(x, y)$, $\psi(x, y)$ of $x$ and $y$: we therefore also know $\phi(x, y) + i\psi(x, y)$, and the question is, whether such a function of $x, y$ (being known when $x + iy$ is known) is to be regarded as a function of $x + iy$; and if not, what is the condition to be satisfied in order that $\phi(x, y) + i\psi(x, y)$ may be a function of $x + iy$. Cauchy at one time considered that the general form was to be regarded as a function of $x + iy$, and he introduced the expression ``\textit{fonction monog\`{e}ne},'' monogenous function, to denote the more restricted form which is the proper function of $x + iy$.

Consider for a moment the above general form, say $x' + iy' = \phi(x, y) + i\psi(x, y)$, where $\phi$, $\psi$ are any real functions of the real variables $(x, y)$; or what is the same thing, assume $x' = \phi(x, y)$, $y' = \psi(x, y)$: if these functions have each or either of them more than one value, we attend only to one value of each of them. We may then as before take $x, y$ to be the coordinates of a point $P$ in a plane $\Pi$, and $x', y'$ to be the coordinates of a point $P'$ in a plane $\Pi'$. If, for any given values of $x, y$, the increments of $\phi(x, y)$, $\psi(x, y)$ corresponding to the indefinitely small real increments $h, k$ of $x, y$ be $Ah + Bk$, $Ch + Dk$, $A, B, C, D$ being functions of $x, y$, then if the new coordinates of $P$ are $x + h$, $y + k$, the new coordinates of $P'$ will be $x' + Ah + Bk$, $y' + Ch + Dk$; or $P$, $P'$ will respectively describe the indefinitely small straight paths at the inclinations $\tan^{-1}\frac{k}{h}$, $\tan^{-1}\frac{Ch + Dk}{Ah + Bk}$ to the axes of $x$, $x'$ respectively; calling these angles $\theta$, $\theta'$, we have therefore $\tan\theta' = \frac{C + D\tan\theta}{A + B\tan\theta}$. And if $x' + iy'$ be $= f(x + iy)$, a function of $x + iy$, the condition to be satisfied is that the increment of $x' + iy'$ shall be proportional to the increment $h + ik$ of $x + iy$, or say that it shall be $= (\lambda + i\mu)(h + ik)$, $\lambda, \mu$ being functions of $x, y$, but independent of $h, k$; we must therefore have $Ah + Bk$, $Ch + Dk$ $= \lambda h - \mu k$, $\mu h + \lambda k$ respectively, that is $A, B, C, D = \lambda, -\mu, \mu, \lambda$ respectively, and the equation for $\tan\theta'$ thus becomes $\tan\theta' = \frac{\mu + \lambda\tan\theta}{\lambda - \mu\tan\theta}$; hence writing $\frac{\mu}{\lambda} = \tan\alpha$, where $\alpha$ is a function of $x, y$, but independent of $h, k$, we have $\tan\theta' = \frac{\tan\alpha + \tan\theta}{1 - \tan\alpha\tan\theta}$; that is, $\theta' = \alpha + \theta$; or for the given points $(x, y)$, $(x', y')$, the path of $P$ being at any inclination whatever $\theta$ to the axis of $x$, the path of $P'$ is at the inclination $\theta + \text{constant angle } \alpha$ to the axis of $x$; also $(\lambda h - \mu k)^{2} + (\mu h + \lambda k)^{2} = (\lambda^{2} + \mu^{2})(h^{2} + k^{2})$, i.e., the lengths of the paths are in a constant ratio.

The condition may be written $\delta(x' + iy') = (\lambda + i\mu)(\delta x + i\delta y)$; or what is the same thing
\[
\frac{dx'}{dx} + i\frac{dy'}{dx} = (\lambda + i\mu)\left(1 + i\frac{dy}{dx}\right),
\]
that is,
\[
\frac{dx'}{dx} = \lambda - \mu\frac{dy}{dx}, \quad \frac{dy'}{dx} = \mu + \lambda\frac{dy}{dx};
\]
consequently
\[
\frac{dx'}{dy} + i\frac{dy'}{dy} = (\lambda + i\mu)\left(\frac{dx}{dy} + i\right),
\]
that is,
\[
\frac{dx'}{dy} = \lambda\frac{dx}{dy} - \mu, \quad \frac{dy'}{dy} = \mu\frac{dx}{dy} + \lambda;
\]
as the analytical conditions in order that $x' + iy'$ may be a function of $x + iy$. They obviously imply $\frac{d^{2}x'}{dx^{2}} + \frac{d^{2}x'}{dy^{2}} = 0$, $\frac{d^{2}y'}{dx^{2}} + \frac{d^{2}y'}{dy^{2}} = 0$; and if $x'$ be a function of $x, y$, satisfying the first of these conditions, then $\frac{dx'}{dx}\,dx + \frac{dx'}{dy}\,dy$ is a complete differential, and $\int\left(\frac{dx'}{dx}\,dx + \frac{dx'}{dy}\,dy\right)$ is a function of $x + iy$.

\subsection*{\S\,20.}

We have, in what just precedes, the ordinary behaviour of a function $\phi(x + iy)$ in the neighbourhood of the value $x + iy$ of the argument or point $x + iy$; or say the behaviour in regard to a point $x + iy$ such that the function is in the neighbourhood of this point a continuous function of $x + iy$ (or that the point is not a point of discontinuity): the correlative definition of continuity will be that the function $\phi(x + iy)$, assumed to have at the given point $x + iy$ a single finite value, is continuous in the neighbourhood of this point, when the point $x + iy$ describing continuously a straight infinitesimal element $h + ik$, the point $\phi(x + iy)$ describes continuously a straight infinitesimal element $(\lambda + i\mu)(h + ik)$: or what is really the same thing, when the function $\phi(x + iy)$ has at the point $x + iy$ a differential coefficient.

\subsection*{\S\,21.}

It would doubtless be possible to give for the continuity of a function $\phi(x + iy)$ a less stringent definition not implying the existence of a differential coefficient; but we have this theory only in regard to the functions of a real variable in memoirs by Riemann, Hankel, du Bois Reymond, Schwarz, Gilbert, Klein, and Darboux. The last-mentioned geometer, in his ``M\'{e}moire sur les fonctions discontinues,'' \textit{Jour.\ de l'\'{E}cole Normale}, t.~iv.\ (1875), pp.~57--112, gives (after Bonnet) the following definition of a continuous function (observe that we are now dealing with real quantities only):

The function $f(x)$ is continuous for the value $x = x_{0}$, when, $h$ and $k$ being positive quantities as small as we please and $\delta$ any positive quantity at pleasure between 0 and 1, we have, for all the values of $\delta$, $f(x_{0} \pm \delta h) - f(x_{0})$ less in absolute magnitude than $\epsilon$; and moreover $f(x)$ is continuous through the interval $x_{0}, x_{1}$ ($x_{1} > x_{0}$, that is, nearer $+\infty$) when $f(x)$ is continuous for every value of $x$ between $x_{0}$ and $x_{1}$, and, $h$ tending to zero through positive values, $f(x_{0} + h)$ and $f(x_{1} - h)$ tend to the limits $f(x_{0})$, $f(x_{1})$ respectively. It is possible, consistently with this definition, to form continuous functions not having in any proper sense a differential coefficient, and having other anomalous properties; thus if $a_{1}, a_{2}, a_{3}, \ldots$ be an infinite series of real positive or negative quantities, such that the series $\Sigma a_{n}$ is absolutely convergent (i.e.\ the sum $\Sigma \pm a_{n}$, each term being made positive, is convergent), then the function $\Sigma a_{n}(\sin n\pi x)^{2}$ is a continuous function actually calculable for any assumed value of $x$: but it is shown in the memoir that, taking $x = \text{any commensurable value } \frac{m}{n}$ whatever, and then writing $x = \frac{m}{n} + h$, $h$ indefinitely small, the increment of the function is of the form $(k + \epsilon)h^{\mu}$, $k$ finite, $\epsilon$ an indefinitely small quantity vanishing with $h$; there is thus no term varying with $h$, nor consequently any differential coefficient. See also Riemann's Memoir, \textit{Ueber die Darstellbarkeit}, \&c.\ (No.~xii.\ in the collected works), already referred to.

\subsection*{\S\,22.}

It was necessary to allude to the foregoing theory of (as they may be termed) infinitely discontinuous functions; but the ordinary and most important functions of analysis are those which are continuous, except for a finite number (or it may be an infinite number) of points of discontinuity. It is to be observed that a point at which the function becomes infinite is \textit{ipso facto} a point of discontinuity; a value of the variable for which the function becomes infinite is, as already mentioned, said to be an ``infinity'' (or a ``pole'') of the function; thus, in the case of a rational function expressed as a fraction in its least terms, if the denominator contains a factor $(x - a)^{m}$, $a$ a real or imaginary value, $m$ a positive integer, then $a$ is said to be an infinity of the $m$th order (and in the particular case $m = 1$, it is said to be a simple infinity). The circular functions $\tan x$, $\sec x$ are instances of a function having an infinite number of simple infinities.

A rational function is a one-valued function, and in regard to a rational function the infinities are the only points of discontinuity; but a one-valued function may have points of discontinuity of a character quite distinct from an infinity: for instance, in the exponential function $\exp\frac{1}{u - a}$, where $a$ is real or imaginary, the value $u (= x + iy) = a$, is a point of discontinuity but not an infinity; taking $u = a + \rho e^{i\alpha}$, where $\rho$ is an indefinitely small real positive quantity, the value of the function is $\exp\frac{1}{\rho e^{i\alpha}} = \exp\frac{1}{\rho}(\cos\alpha - i\sin\alpha)$, which is indefinitely large or indefinitely small according as $\cos\alpha$ is positive or negative, and in the separating case $\cos\alpha = 0$, and therefore $\sin\alpha = \pm 1$, it is $= \cos\frac{1}{\rho} \mp i\sin\frac{1}{\rho}$, which is indeterminate. If, instead of $\exp\frac{1}{u - a}$ we consider a linear function
\[
\{A + B\exp\tfrac{1}{u - a}\} \cdot \{C + D\exp\tfrac{1}{u - a}\},
\]
then writing as before $u = a + \rho e^{i\alpha}$, the value is $= A + C$, or $= B + D$, according as $\cos\alpha$ is negative or positive. As regards the theory of one-valued functions in general, the memoir by Weierstrass, ``Zur Theorie der eindeutigen analytischen Functionen,'' \textit{Berl.\ Abh.}\ 1876, pp.~11--60, may be referred to.

\subsection*{\S\,23.}

A one-valued function \textit{ex vi termini} cannot have a point of discontinuity of the kind next referred to; if the representative point $P$, moving in any manner whatever, returns to its original position, the corresponding point $P'$ cannot but return to its original position. But consider a many-valued function, say an $n$-valued function $x' + iy'$, of $x + iy$: to each position of $P$ there correspond $n$ positions, in general all of them different, of $P'$. But the point $P$ may be such that (to take the most simple case) two of the corresponding points $P'$ coincide with each other, say such a position of $P$ is at $V$, then (using for greater distinctness a different letter $W$ instead of $V$) corresponding thereto we have two coincident points $(W)$, and $n - 2$ other points $W$; $V$ is then a branch-point (\textit{Verzweigungspunkt}). Taking for $P$ any point which is not a branch-point, then in the neighbourhood of this value each of the $n$ functions $x' + iy'$ is a continuous function of $x + iy$, and by what precedes, if $P$ describing an infinitely small closed curve (or oval) return to its original position, then each of the corresponding points $P'$ describing a corresponding indefinitely small oval will return to its original position. But imagine the oval described by $P$ to be gradually enlarged, so that it comes to pass through a branch-point $V$; the ovals described by two of the corresponding points $P'$ will gradually approach each other, and will come to unite at the point $(W)$, each oval then sharpening itself out so that the two form together a figure of eight. And if we imagine the oval described by $P$ to be still further enlarged so as to include within it the point $V$, then the figure of eight, losing the crossing, will be at first an hour-glass form, or twice-indented oval, and ultimately in form an ordinary oval, but having the character of a twofold oval; i.e.\ to the oval described by $P$ (and which surrounds the branch-point $V$) there will correspond this twofold oval, and $n - 2$ onefold ovals, in such wise that to a given position of $P$ on its oval there correspond two points, say $P_{1}'$, $P_{2}'$, on the twofold oval, and $n - 2$ points $P_{3}', \ldots, P_{n}'$, each on its own onefold oval. And then as $P$ describing its oval returns to its original position, the point $P_{1}'$ describing a portion only of the twofold oval, will pass to the original position of $P_{2}'$, while the point $P_{2}'$ describing the remaining portion of the twofold oval will pass to the original position of $P_{1}'$; the other points $P_{3}', \ldots, P_{n}'$, describing each of them its own onefold oval, will return each of them to its original position. And it is easy to understand how, when the oval described by $P$ surrounds two or more of the branch-points $V$, the corresponding curves for $P'$ may be a system of manifold ovals, such that the sum of the manifoldness is always $= n$, and to conceive in a general way the behaviour of the corresponding points $P$ and $P'$.

Writing for a moment $x + iy = u$, $x' + iy' = v$, the branch-points are the points of contact of parallel tangents to the curve $\phi(u, v) = 0$, a line through a cusp (but not a line through a node), being reckoned as a tangent: that is, if this be a curve of the order $n$ and class $m$, with $\delta$ nodes and $\kappa$ cusps, the number of branch-points is $= m + \kappa$, that is, it is $= n^{2} - 2\delta - 2\kappa$, or if $p = \frac{1}{2}(n - 1)(n - 2) - \delta - \kappa$, be the deficiency, then the number is $= 2n - 2 + 2p$.

To illustrate the theory of the $n$-valued algebraical function $x' + iy'$ of the complex variable $x + iy$, Riemann introduces the notion of a surface composed of $n$ coincident planes or sheets, such that the transition from one sheet to another is made at the branch-points, and that the $n$ sheets form together a multiply-connected surface, which can by cross-cuts (\textit{Querschnitte}) be converted into a simply-connected surface: the $n$-valued function $x' + iy'$ becomes thus a one-valued function of $x + iy$, considered as belonging to a point on some determinate sheet of the surface: and upon such considerations he founds the whole theory of the functions which arise from the integration of the differential expressions depending on the $n$-valued algebraical function (that is, any irrational algebraical function whatever) of the independent variable, establishing as part of the investigation the theory of the multiple $\vartheta$-functions. But it would be difficult to give a further account of these investigations.

\smallskip
\begin{center}\textbf{The Calculus of Functions.}\end{center}

\subsection*{\S\,24.}

The so-called Calculus of Functions, as considered chiefly by Herschel and Babbage and De Morgan, is not so much a theory of functions as a theory of the solution of functional equations; or, as perhaps should rather be said, the solution of functional equations by means of known functions, or symbols,---the epithet known being here used in reference to the actual state of analysis. Thus for a functional equation $\phi x + \phi y = \phi(xy)$, taking the logarithm as a known function, the solution is $\phi x = c\log x$; or if the logarithm is not taken to be a known function, then a solution may be obtained by means of the sign of integration $\phi x = c\int\frac{dx}{x}$; but the establishment of the properties of the function logarithm (assumed to be previously unknown) would not be considered as coming within the theory. A class of equations specially considered is where $\alpha x$, $\beta x$, \ldots\ being given functions of $x$, the unknown function $\phi$ is to be determined by means of a given relation between $x$, $\phi x$, $\phi\alpha x$, $\phi\beta x$, \ldots; in particular the given relation may be between $x$, $\phi x$, $\phi\alpha x$; this can be at once reduced to equations of finite differences; for writing $x = u_{n}$, $\alpha x = u_{n+1}$, we have $u_{n+1} = \alpha u_{n}$, giving $u_{n}$, and therefore also $x$, each of them as a function of $n$; and then writing $\phi x = v_{n}$, $\phi\alpha x$ will be the same function of $n + 1$, $= v_{n+1}$, and the given relation is again an equation of finite differences in $v_{n+1}$, $v_{n}$, and $n$; we have thus $v_{n} = \phi x$ as a function of $n$, that is, of $x$. As regards the equation $u_{n+1} = \alpha u_{n}$, considered in itself apart from what precedes, observe that this is satisfied by writing $u_{n} = \alpha^{n}(x)$, or the question of solving this equation of finite differences is, in fact, identical with that of finding the $n$th function $\alpha^{n}(x)$, where $\alpha(x)$ is a given function of $x$. It of course depends on the form of $\alpha(x)$ whether this question admits of solution in any proper sense: thus, for a function such as $\log x$, the $n$th logarithm is expressible in its original form $\log^{n}x$ ($= \log\log\ldots x$), and not in any other form. But there are forms, for instance $\alpha x = \frac{A + Bx}{C + Dx}$, where the $n$th function $\alpha^{n}x$ is a function of the like form $\alpha^{n}x = \frac{A_{n} + B_{n}x}{C_{n} + D_{n}x}$, in which the actual value can be expressed as a function of $n$; if $\alpha$ be such a form, then $\phi\alpha\phi^{-1}$, whatever $\phi$ may be, is a like form, for we obviously have $(\phi\alpha\phi^{-1})^{n} = \phi\alpha^{n}\phi^{-1}$. The determination of the $n$th function is, in fact, a leading question in the calculus of functions.

It is to be observed that considering the case of two variables, if for instance $\phi(x, y)$ denote a given function of $x, y$, the notation $\phi^{n}(x, y)$ is altogether meaningless: in order to generalize the question, we require an extended notation wherein a single functional symbol is used to denote two functions of the two variables. Thus $\phi(x, y) = \alpha(x, y)$, $\beta(x, y)$, $\alpha$ and $\beta$ given functions; writing for shortness $x_{1} = \alpha(x, y)$, $y_{1} = \beta(x, y)$, then $\phi^{2}(x, y)$ will denote $\phi(x_{1}, y_{1})$, that is, two functions $\alpha(x_{1}, y_{1})$, $\beta(x_{1}, y_{1})$, say these are $x_{2}$, $y_{2}$; $\phi^{3}(x, y)$ will denote $\phi(x_{2}, y_{2})$, and so on, so that $\phi^{n}(x, y)$ will have a determinate meaning. And the like is obviously the case in regard to any number of variables, the single functional symbol denoting in each case a set of functions equal in number to the variables.

\newpage
% ============================================================
\section*{\S\,788. GALOIS.}
\addcontentsline{toc}{section}{788. Galois}
% ============================================================

\noindent[From the \textit{Encyclop{\ae}dia Britannica}, Ninth Edition, vol.~x (1879), p.~48.]

\smallskip

\textsc{Galois, Evariste} (1811--1832), an eminently original and profound French mathematician, born 26th October 1811, killed in a duel May 1832. A necrological notice by his friend M.\ Auguste Chevalier appeared in the \textit{Revue Encyclop\'{e}dique}, September 1832, p.~744; and his collected works are published, \textit{Liouville}, t.~xi.\ (1846), pp.~381--444, about fifty of these pages being occupied by researches on the resolubility of algebraic equations by radicals. But these researches, crowning as it were the previous labours of Lagrange, Gauss, and Abel, have in a signal manner advanced the theory, and it is not too much to say that they are the foundation of all that has since been done, or is doing, in the subject. The fundamental notion consists in the establishment of a group of permutations of the roots of an equation, such that every function of the roots invariable by the substitutions of the group is rationally known, and reciprocally that every rationally determinable function of the roots is invariable by the substitutions of the group; some further explanation of the theorem, and in connexion with it an explanation of the notion of an adjoint radical, is given under Equation, No.~32, [786]. As part of the theory (but the investigation has a very high independent value as regards the Theory of Numbers, to which it properly belongs), Galois introduces the notion of the imaginary roots of an irreducible congruence of a degree superior to unity; i.e., such a congruence, $F(x) \equiv 0$ (mod.\ a prime number $p$), has no integer root; but what is done is to introduce a quantity $i$ subjected to the condition of verifying the congruence in question, $F(i) \equiv 0$ (mod.\ $p$), which quantity $i$ is an imaginary of an entirely new kind, occupying in the theory of numbers a position analogous to that of $\sqrt{-1}$ in algebra.

\newpage
% ============================================================
\section*{\S\,789. GAUSS.}
\addcontentsline{toc}{section}{789. Gauss}
% ============================================================

\noindent[From the \textit{Encyclop{\ae}dia Britannica}, Ninth Edition, vol.~x (1879), p.~116.]

\smallskip

\textsc{Gauss, Carl Friedrich} (1777--1855), an eminent German mathematician, was born of humble parents at Brunswick, April 23, 1777, and was indebted for a liberal education to the notice which his talents procured him from the reigning duke. His name became widely known by the publication, in his twenty-fifth year (1801), of the \textit{Disquisitiones Arithmetic\ae}. In 1807 he was appointed director of the G\"{o}ttingen observatory, an office which he retained to his death: it is said that he never slept away from under the roof of his observatory, except on one occasion, when he accepted an invitation from Humboldt to attend a meeting of natural philosophers at Berlin. In 1809 he published at Hamburg his \textit{Theoria Motus Corporum Coelestium}, a work which gave a powerful impulse to the true methods of astronomical observation; and his astronomical workings, observations, calculations of orbits of planets and comets, \&c., are very numerous and valuable. He continued his labours in the theory of numbers and other analytical subjects, and communicated a long series of memoirs to the Royal Society of Sciences at G\"{o}ttingen. His first memoir on the theory of magnetism, \textit{Intensitas vis magnetic\ae} terrestris ad mensuram absolutam revocata, was published in 1833, and he shortly afterwards proceeded, in conjunction with Professor Wilhelm Weber, to invent new apparatus for observing the earth's magnetism and its changes; the instruments devised by them were the declination instrument and the bifilar magnetometer. With Weber's assistance he erected in 1833 at G\"{o}ttingen a magnetic observatory free from iron (as Humboldt and Arago had previously done on a smaller scale), where he made magnetic observations, and from this same observatory he sent telegraphic signals to the neighbouring town, thus showing the practicability of an electromagnetic telegraph. He further instituted an association (Magnetische Verein), composed at first almost entirely of Germans, whose continuous observations on fixed term-days extended from Holland to Sicily. The volumes of their publication, \textit{Resultate aus den Beobachtungen des magnetischen Vereins}, extend from 1836 to 1839; and in those for 1838 and 1839 are contained the two important memoirs by Gauss, \textit{Allgemeine Theorie des Erdmagnetismus}, and the \textit{Allgemeine Lehrs\"{a}tze}---on the theory of forces attracting according to the inverse square of the distance. The instruments and methods thus due to him are substantially those employed in the magnetic observatories throughout the world. He co-operated in the Danish and Hanoverian measurements of an arc and trigonometrical operations (1821--48), and wrote (1843, 1846) the two memoirs \textit{Ueber Gegenst\"{a}nde der h\"{o}hern Geod\"{a}sie}. He co-operated in the Danish and Hanoverian measurements of an arc and trigonometrical operations (1821--48), and wrote (1843, 1846) the two memoirs \textit{Ueber Gegenst\"{a}nde der h\"{o}hern Geod\"{a}sie}. Connected with observations in general we have (1812--26) the memoir \textit{Theoria combinationis observationum erroribus minimis obnoxi\ae}, with a second part and a supplement. Another memoir of applied mathematics is the \textit{Dioptrische Untersuchungen}, 1840. Gauss was well versed in general literature and the chief languages of modern Europe, and was a member of nearly all the leading scientific societies in Europe. He died at G\"{o}ttingen early in the spring of 1855. The centenary of his birth was celebrated (1877) at his native place, Brunswick.

Gauss's collected works have been recently published by the Royal Society of G\"{o}ttingen, in 7 vols.\ 4to, Gott., 1863--71, edited by E.\ J.\ Schering,---(1) the \textit{Disquisitiones Arithmetic\ae}, (2) Theory of Numbers, (3) Analysis, (4) Geometry and Method of Least Squares, (5) Mathematical Physics, (6) Astronomy, and (7) the \textit{Theoria Motus Corporum Coelestium}. They include, besides his various works and memoirs, notices by him of many of these, and of works of other authors in the G\"{o}ttingische gelehrte Anzeigen, and a considerable amount of previously unpublished matter, \textit{Nachlass}. Of the memoirs in pure mathematics, comprised for the most part in vols.~II., III., and IV.\ (but to these must be added those on Attractions in vol.~v.), it may be safely said there is not one which has not signally contributed to the progress of the branch of mathematics to which it belongs, or which would not require to be carefully analysed in a history of the subject. Running through these volumes in order, we have in the second the memoir, \textit{Summatio quarundam, serierum singularium}, the memoirs on the theory of biquadratic residues, in which the notion of complex numbers of the form $a + bi$ was first introduced into the theory of numbers: and included in the \textit{Nachlass} are some valuable tables. That for the conversion of a fraction into decimals (giving the complete period for all the prime numbers up to 997) is a specimen of the extraordinary love which Gauss had for long arithmetical calculations; and the amount of work gone through in the construction of the table of the number of the classes of binary quadratic forms must also have been tremendous. In vol.~iii.\ we have memoirs relating to the proof of the theorem that every numerical equation has a real or imaginary root, the memoirs on the Hypergeometric Series, that on Interpolation, and the memoir \textit{Determinatio Attractionis} in which a planetary mass is considered as distributed over its orbit according to the time in which each portion of the orbit is described, and the question (having an implied reference to the theory of secular perturbations) is to find the attraction of such a ring. In the solution the value of an elliptic function is found by means of the arithmetico-geometrical mean. The \textit{Nachlass} contains further researches on this subject, and also researches (unfortunately very fragmentary) on the lemniscate-function, \&c., showing that Gauss was, even before 1800, in possession of many of the discoveries which have made the names of Abel and Jacobi illustrious. In vol.~iv.\ we have the memoir \textit{Allgemeine Aufl\"{o}sung\ldots}, on the graphical representation of one surface upon another, and the \textit{Disquisitiones generales circa superficies curvas}. And in vol.~v.\ we have a memoir On the Attraction of Homogeneous Ellipsoids, and the already mentioned memoir \textit{Allgemeine Lehrs\"{a}tze\ldots}, on the theory of forces attracting according to the inverse square of the distance.

\newpage
% ============================================================
\section*{\S\,790. GEOMETRY (ANALYTICAL).}
\addcontentsline{toc}{section}{790. Geometry (Analytical)}
% ============================================================

\noindent[From the \textit{Encyclop{\ae}dia Britannica}, Ninth Edition, vol.~x (1879), pp.~408--420.]

\smallskip

This will be here treated as a method. The science is Geometry; and it would be possible, analytically, or by the method of coordinates, to develope the truths of geometry in a systematic course. But it is proposed not in any way to attempt this, but simply to explain the method, giving such examples, interesting (it may be) in themselves, as are suitable for showing how the method is employed in the demonstration and solution of theorems or problems.

Geometry is one-, two-, or three-dimensional, or, what is the same thing, it is lineal, plane, or solid, according as the space dealt with is the line, the plane, or ordinary (three-dimensional) space. No more general view of the subject need here be taken: but in a certain sense one-dimensional geometry does not exist, inasmuch as the geometrical constructions for points in a line can only be performed by travelling out of the line into other parts of a plane which contains it, and conformably to the usual practice Analytical Geometry will be treated under the two divisions, Plane and Solid.

It is proposed to consider Cartesian coordinates almost exclusively; for the proper development of the science homogeneous coordinates (three and four in plane and solid geometry respectively) are required: and it is moreover necessary to have the correlative line- and plane-coordinates: and in solid geometry to have the six coordinates of the line. The most comprehensive English works are those of Dr Salmon, \textit{Conic Sections} (5th edition, 1869), \textit{Higher Plane Curves} (2nd edition, 1873), and \textit{Geometry of Three Dimensions} (3rd edition, 1874): we have also, on plane geometry, Clebsch's \textit{Vorlesungen \"{u}ber Geometrie}, posthumous, edited by Dr F.\ Lindemann, Leipsic, 1875, not yet complete.

\smallskip
\begin{center}\textbf{I.\ Plane Analytical Geometry (\S\S 1--25).}\end{center}

\subsection*{\S\,1.}

It is assumed that the points, lines, and figures considered exist in one and the same plane, which plane, therefore, need not be in any way referred to. The position of a point is determined by means of its (Cartesian) coordinates: i.e.\ as explained under the article Curve, [785], we take the two lines $x'Ox$ and $y'Oy$, called the \textit{axes} of $x$ and $y$ respectively, intersecting in a point $O$ called the \textit{origin}, and determine the position of any other point $P$ by means of its coordinates $x = OM$ (or $NP$), and $y = MP$ (or $ON$). The two axes are usually (as in fig.~1) at right angles to each other, and the lines $PM$, $PN$ are then at right angles to the axes of $x$ and $y$ respectively. Assuming a scale at pleasure, the coordinates $x$, $y$ of a point have numerical values.

It is necessary to attend to the signs: $x$ has opposite signs according as the point is on one side or the other of the axis of $y$, and similarly $y$ has opposite signs according as the point is on the one side or the other of the axis of $x$. Using the letters N., E., S., W.\ as in a map, and considering the plane as divided into four \textit{quadrants} by the axes, the signs are usually taken to be
\begin{center}
\begin{tabular}{c|c|c}
$x$ & $y$ & quad. \\ \hline
$+$ & $+$ & N.E. \\
$+$ & $-$ & S.E. \\
$-$ & $+$ & N.W. \\
$-$ & $-$ & S.W.
\end{tabular}
\end{center}
A point is said to have the coordinates $(a, b)$, and is referred to as \textit{the point $(a, b)$}, when its coordinates are $x = a$, $y = b$; the coordinates $x$, $y$ of a variable point, or of a point which is for the time being regarded as variable, are said to be \textit{current coordinates}.

\subsection*{\S\,2.}

It is sometimes convenient to use oblique coordinates; the only difference is that the axes are not at right angles to each other; the lines $PM$, $PN$ are drawn parallel to the axes of $y$ and $x$ respectively, and the figure $OMPN$ is thus a parallelogram. But in all that follows, the Cartesian coordinates are taken to be rectangular; polar coordinates and other systems will be briefly referred to in the sequel.

\subsection*{\S\,3.}

If the coordinates $(x, y)$ of a point are not given, but only a relation between them $f(x, y) = 0$, then we have a \textit{curve}. For, if we consider $x$ as a real quantity varying continuously from $-\infty$ to $+\infty$, then, for any given value of $x$, $y$ has a value or values. If these are all imaginary, there is not any real point; but if one or more of them be real, we have a real point or points, which (as the assumed value of $x$ varies continuously) varies or vary continuously therewith; and the locus of all these real points is \textit{a curve}. The equation completely defines the curve: to trace the curve directly from the equation, nothing else being known, we obtain as above a series of points sufficiently near to each other, and draw the curve through them. For instance, let this be done in a simple case. Suppose $y = 2x - 1$; it is quite easy to obtain and lay down a series of points as near to each other as we please, and the application of a ruler would show that these were in a line: that the curve is a line depends upon something more than the equation itself, viz.\ the theorem that every equation of the form $y = ax + b$ represents a line: supposing this known, it will be at once understood how the process of tracing the curve may be abbreviated; we have $x = 0$, $y = -1$, and $x = \frac{1}{2}$, $y = 0$; the curve is thus the line passing through these two points. But in the foregoing example the notion of a line is taken to be a known one, and such notion of a line does in fact precede the consideration of any equation of a curve whatever, since the notion of the coordinates themselves rests upon that of a line. In other cases it may very well be that the equation is the definition of the curve; the points laid down, although (as finite in number) they do not actually determine the curve, determine it to any degree of accuracy; and the equation thus enables us to construct the curve.

A curve may be determined in another way; viz.\ the coordinates $x, y$ may be given each of them as a function of the same variable \textit{parameter}; $x, y = f(\theta)$, $\phi(\theta)$ respectively. Here, giving to $\theta$ any number of values in succession, these equations determine the values of $x, y$, that is, the positions of a series of points on the curve. The ordinary form $y = \phi(x)$, where $y$ is given explicitly as a function of $x$, is a particular case of each of the other two forms: we have $f(x, y) = y - \phi(x) = 0$; and $x = \theta$, $y = \phi(\theta)$.

\subsection*{\S\,4.}

As remarked under Curve, [785], it is a useful exercise to trace a considerable number of curves, first taking equations which are purely numerical, and then equations which contain literal constants (representing numbers): the equations most easily dealt with are those wherein one coordinate is given as an explicit function of the other, say $y = \phi(x)$ as above. A few examples are here given, with such explanations as seem proper.

(i) $y = 2x - 1$, as before; it is at once seen that this is a line; and taking it to be so, any two points, for instance, $(0, -1)$ and $(\frac{1}{2}, 0)$, determine the line.

(ii) $y = x^{2}$. The equation shows that $x$ may be positive or negative, but that $y$ is always positive, and has the same values for equal positive and negative values of $x$: the curve passes through the origin, and through the points $(\pm 1, 1)$. It is already known that the curve lies wholly above the axis of $x$. To find its form in the neighbourhood of the origin, give $x$ a small value, $x = \pm 0.1$ or $\pm 0.01$, then $y$ is very much smaller, $= 0.01$ and $0.0001$ in the two cases respectively: this shows that the curve touches the axis of $x$ at the origin. Moreover, $x$ may be as large as we please, but when it is large, $y$ is much larger; for instance, $x = 10$, $y = 100$. The curve is a \textit{parabola} (fig.~2).

(iii) $y = x^{3}$. Here $x$ being positive $y$ is positive, but $x$ being negative $y$ is also negative: the curve passes through the origin, and also through the points $(1, 1)$ and $(-1, -1)$. Moreover, when $x$ is small, $= 0.1$ for example, then not only is $y$, $= 0.001$, very much smaller than $x$, but it is also very much smaller than $y$ was for the last-mentioned curve $y = x^{2}$, that is, in the neighbourhood of the origin the present curve approaches more closely the axis of $x$. The axis of $x$ is a tangent at the origin, but it is a tangent of a peculiar kind (a stationary or inflexional tangent), cutting the curve at the origin, which is an \textit{inflexion}. The curve is the \textit{cubical parabola} (fig.~3).

(iv) $y^{2} = x - 1 \cdot x - 3 \cdot x - 4$. Here $y = 0$ for $x = 1$, $= 3$, $= 4$. Whenever $x - 1 \cdot x - 3 \cdot x - 4$ is positive, $y$ has two equal and opposite values: but when $x - 1 \cdot x - 3 \cdot x - 4$ is negative, then $y$ is imaginary. In particular, for $x$ less than 1, or between 3 and 4, $y$ is imaginary, but for $x$ between 1 and 3, or greater than 4, $y$ has two values. It is clear that for $x$ somewhere between 1 and 3, $y$ will attain a maximum: the values of $x$ and $y$ may be found approximately by trial. The curve will consist of an oval and infinite branch, and it is easy to see that, as shown in fig.~4, the curve where it cuts the axis of $x$ cuts it at right angles. It may be further remarked that, as $x$ increases from 4, the value of $y$ will increase more and more rapidly: for instance, $x = 6$, $y^{2} = 8$, $x = 10$, $y^{2} = 378$, \&c., and it is easy to see that this implies that the curve has on the infinite branch two inflexions as shown.

(v) $y^{2} = x - c \cdot x - b \cdot x - a$, where $a > b > c$ (that is, $a$ nearer to $+\infty$, $c$ to $-\infty$). The curve has the same general form as in the last figure, the oval extending between the limits $x = c$, $x = b$, the infinite branch commencing at the point $x = a$.

(vi) $y^{2} = (x - c)^{2}(x - a)$. Suppose that in the last-mentioned curve, $y^{2} = x - c \cdot x - b \cdot x - a$, $b$ gradually diminishes, and becomes ultimately $= c$. The infinite branch (see fig.~5) changes its form, but not in a very marked manner, and it retains the two inflexions. The oval lies always between the values $x = c$, $x = b$, and therefore its length continually diminishes: it is easy to see that its breadth will also continually diminish: ultimately it shrinks up into a mere point. The curve has thus a \textit{conjugate or isolated point}, or \textit{acnode}. For a direct verification observe that $x = c$, $y = 0$, so that $(c, 0)$ is a point of the curve, but if $x$ is either less than $c$, or between $c$ and $a$, $y^{2}$ is negative, and $y$ is imaginary.

(vii) $y^{2} = (x - c)(x - a)^{2}$. If in the same curve $b$ gradually increases and becomes ultimately $= a$, the oval and the infinite branch change each of them its form, the oval extending always between the values $x = c$, $x = b$, and thus continually approaching the infinite branch, which begins at $x = a$. The consideration of a few numerical examples, with careful drawing, would show that the oval and the infinite branch as they approach sharpen out each towards the other, the two inflexions on the infinite branch coming always nearer to the point $(a, 0)$,---so that finally, when $b$ becomes $= a$, the curve has the form shown in fig.~6, there being now a double point or node (crunode) at $A$, and the inflexions on the infinite branch having disappeared.

\vspace{\baselineskip}
\begin{center}
\rule{0.5\textwidth}{0.4pt}\\[4pt]
\textit{End of transcription---pages 501--550.}
\end{center}

\end{document}
